\documentclass[../main.tex]{subfiles}
\begin{document}

\chapter{GPUシミュレーション}
\lessoninfo{
  video={M8，レッスン 1--4（GPU Cycle Level Simulation，Parts 1 and 2；Analytical Models of GPUs；Accelerating GPU Simulation）},
  sources={Liu ら \cite{liu2023photon}（Photon，指定）；Williams ら \cite{williams2009roofline}（Roofline，指定）；Hong と Kim \cite{hong2009analytical}（MWP/CWP モデル）；Huang ら \cite{huang2014gpumech}（GPUMech）；Karkhanis と Smith \cite{karkhanis2004firstorder}（一次近似モデル）；シミュレータ \cite{bakhoda2009gpgpusim,khairy2020accelsim,sun2019mgpusim,gthparch2026macsim,power2015gem5gpu}；Sherwood ら \cite{sherwood2002simpoint}（SimPoint）；NVIDIA GEMM ガイド \cite{nvidia2026gemmguide}},
  textbook={Eeckhout \cite{eeckhout2010perfeval}（任意）；K12 \cite{kim2012gpgpu} ch.~2 Performance Principles，§4.1.1 Execution Time Modeling，§4.2.4 Roofline Model，§4.2.6 Other GPGPU Performance Modeling Techniques；BPG \cite{nvidia2026cudabestpractices} §10.2.1 Coalesced Access to Global Memory},
  keywords={性能モデリング，サイクルレベルシミュレーション，実行駆動シミュレーションとトレース駆動シミュレーション，キューベースモデル，セクタキャッシュ，解析モデル，ルーフラインモデル，演算強度，リッジポイント，CPI スタック，インターバル解析，MWP と CWP，サブコア，イベント駆動シミュレーション，サンプリングシミュレーション，SimPoint，データ駆動モデル},
}

これまでの章で扱った機構（第6章の SIMT スタックと動的ワープ形成，
レジスタファイルキャッシュ，第7章のワープスケジューラ）は，いずれも
シリコン上に存在する前に提案され，評価されたものである．本章は，
そうした評価を可能にする道具を扱う．まず，キューとクロックから GPU の
サイクルレベルシミュレータを組み立て，次に GPU に固有のもの，すなわち
ワープ，ダイバージェンス，コアレッシング，セクタキャッシュを加える．
続いて，何もシミュレーションせずに少数のパラメータから性能を予測する
解析モデルに移る．GPU の一次近似設計，ルーフラインモデル，CPI
スタック，インターバル解析である．最後に，詳細なシミュレーションが
なぜこれほど遅いのか，それに対して何ができるのかを問う．並列
シミュレータとイベント駆動シミュレータ，より単純なモデル，
サンプリング，より小さなワークロード，データから学習したモデルである．

\section{サイクルレベルシミュレーション}
\label{sec:l08-cycle-level}

プロセッサを作る前に，設計者はそれが実際のプログラムをどれだけ速く
実行するか，そして設計上の各選択がその速度をどう変えるかを予測
しなければならない．\source{M8，サイクルレベルシミュレーション（パート~1）；Eeckhout
\cite{eeckhout2010perfeval}．}本節では，そのための手法を列挙し，
サイクルレベルシミュレーションとその2つの方式を定義し，パイプラインの
シミュレータがクロックで進められるキューとしてどう構成されるかを示す．

\subsection{6つの手法}

\begin{definition}[性能モデリング]
  \label{def:l08-performance-modeling}
  プロセッサの設計の性能，あるいはある設計上でのプログラムの性能を，
  完成したハードウェアで測定せずに予測することを
  \term[せいのうもでりんぐ]{性能モデリング}[performance modeling]という．
  そのような予測を出せるだけの詳しさでハードウェアの振る舞いを模倣する
  プログラムが\term[しみゅれーた]{シミュレータ}[simulator]である．
\end{definition}

一般に使われる手法は6つある．
\begin{enumerate}
  \item \emph{サイクルレベルシミュレーション}：マシンの各部分が各サイクルに
    何をするかをモデル化する（本節と次節）．
  \item \emph{イベント駆動シミュレーション}：何かが変化する時点だけを
    モデル化する（\cref{sec:l08-accelerating}）．
  \item \emph{解析モデル}：少数のパラメータによる方程式
    （\cref{sec:l08-analytical}）．
  \item \emph{サンプリングに基づく手法}：プログラムの選んだ部分だけを
    詳細にシミュレーションする（\cref{sec:l08-accelerating}）．
  \item \emph{データ駆動モデル，統計モデル，機械学習モデル}：
    測定したデータやシミュレーションで得たデータにモデルを当てはめる
    （\cref{sec:l08-accelerating}）．
  \item
    \term[えふぴーじーえーえみゅれーしょん]{FPGAエミュレーション}[FPGA emulation]：
    設計またはそのモデルを FPGA（field-programmable gate array）上に
    実装する．ソフトウェアシミュレータより桁違いに速く動くが，
    設計の手間がはるかに大きい．\margin{FPGA エミュレーションは
    サイクルレベルシミュレーションと同じだけ詳細で，速いだけである．
    FPGA ベースのシミュレータは\enquote{fast and accurate, but they take
    significant effort to develop and use}とされる
    \cite{sanchez2013zsim}．本章ではこれ以上扱わない．}
\end{enumerate}

\subsection{サイクルレベルの手法}

\begin{definition}[サイクルレベルシミュレーション]
  \label{def:l08-cycle-level}
  \term[さいくるれべるしみゅれーしょん]{サイクルレベルシミュレーション}[cycle-level simulation]
  では，シミュレータはグローバルなクロックを持ち，シミュレーションする
  各サイクルで，各ユニットにおけるそのサイクルの事象をモデル化する．
  どの命令がフェッチされ，どれがデコードされ，どれがスケジューラを
  出て，どの要求がキャッシュに届くか，である．
\end{definition}

サイクルレベルシミュレータはアーキテクチャ研究の標準的な道具であり，
特に測定すべきハードウェアがまだない設計の初期に使われる．クロックが
1つというのは単純化である．実際のチップは，コア，メモリ，オンチップ
ネットワークを異なる周波数で動かす．設計のうち1つのクロックで駆動される
部分を\term[くろっくどめいん]{クロックドメイン}[clock domain]と呼び，
シミュレータはドメイン（コア，メモリ，ネットワークオンチップ）ごとに
1つのクロックを持ち，それぞれを固有の速さで進めることができる．\margin{NVIDIA
GPU では SM のクロックとメモリのクロックは独立している．
\texttt{nvidia-smi -q -d CLOCK} は\enquote{Graphics}，\enquote{SM}，
\enquote{Memory}の各クロックを別々に報告する．}

\subsection{実行駆動シミュレーションとトレース駆動シミュレーション}

サイクルレベルシミュレータはタイミングをモデル化する．それに加えて，
\emph{どの}命令が，どのアドレスとどの分岐結果で実行されるかも知る
必要があり，それを得る方法は2つある．

\begin{definition}[実行駆動シミュレーションとトレース駆動シミュレーション]
  \label{def:l08-exec-trace}
  \term[じっこうくどうしみゅれーしょん]{実行駆動シミュレーション}[execution-driven simulation]
  では，シミュレータ自身がプログラムの命令を実行し，シミュレーションの
  進行に合わせてその結果を計算する．
  \term[とれーすくどうしみゅれーしょん]{トレース駆動シミュレーション}[trace-driven simulation]
  では，実際の実行または機能的な実行から事前に収集した命令（および
  そのアドレスと分岐結果）の記録である\term[とれーす]{トレース}[trace]
  を再生する．こうして実行とタイミングが切り離される．
\end{definition}

実行駆動シミュレータでも，命令を\emph{いつ}機能的に実行するかは選ぶ
必要がある．命令のフェッチ時に実行すれば（\emph{execute-at-fetch}）
機能モデルは単純に保たれ，タイミングモデルが実行ステージに達したときに
実行すれば（\emph{execute-at-execute}），値がシミュレーション上の事象の
順序に依存できる．\caution{2つの選択は通常同じ結果を与える．違いが出るのは
シミュレーション上のタイミングが値を変える場合で，たとえば2つの
スレッドがロックを奪い合うときである．execute-at-fetch では勝者は
フェッチ順で決まり，execute-at-execute ではシミュレーション上の実行順で
決まる．}2つの方式の得失は次のとおりである．
\begin{description}
  \item[実行駆動．]  トレースを保存する必要がなく，実行そのものに依存する
    振る舞いを忠実にモデル化できる．どのスレッドが先にロックを獲得するか，
    スレッドがバリアでどれだけ待つか，データ依存の分岐が異なる
    スケジュールのもとでどの経路をとるか，である．
  \item[トレース駆動．]  シミュレータは機能モデルを必要としないので，
    より単純で軽く，開発しやすい．キャッシュだけを調べるなら，メモリ
    アドレス（メモリトレース）があれば足りる．しかしトレースは，記録時の
    事象の順序を固定する．ロックの獲得，バリア，そのほか実行時の
    タイミングに依存するものは，シミュレーション上のマシンで変化できない．
\end{description}
\tip{経験則：問いが\enquote{この設計はどれだけ速いか}であればトレースで
足りる．設計がプログラムの\emph{振る舞いそのもの}（同期やスレッドの
たどる経路）を変えうるなら，シミュレーションは実行駆動でなければならない．}

\subsection{キューベースのモデル化}

第4章のパイプラインは，キューで結ばれたステージとしてモデル化するのが
最も簡単である（\cref{fig:l08-queue-model}）．

\begin{definition}[キューベースモデル]
  \label{def:l08-queue-model}
  \term[きゅーべーすもでる]{キューベースモデル}[queue-based model]は，
  パイプラインをキューで結ばれたステージの集合として表す．各キューは
  \term{FIFO}[FIFO]（first in, first out）であり，命令は入った順に
  出ていく．例外はスケジューラで，\emph{準備のできた}任意の命令を
  選ぶので，FIFO ではなくプールである．各接続は2つのパラメータで
  記述される．キューの\emph{深さ}は命令がそのステージで過ごすサイクル数で
  あり，\term[ぱいぷらいんはば]{パイプライン幅}[pipeline width]は
  1サイクルにあるキューから次のキューへ移動できる命令数（フェッチ，
  発行，実行の帯域）である．
\end{definition}

\begin{figure}[tb]
  \centering
  % Queue-based model of a pipeline (Lesson 8).  The stages fetch, decode,
% schedule, execute and write-back are connected by FIFO queues
% (FE, DE and EX queue); the scheduler is drawn as a pool from which
% ready instructions are picked, not as a FIFO.  Execute talks to the
% memory model; completion wakes up dependent instructions.
% Redraws the outline figure m8l1_queue.
% Usage: % Queue-based model of a pipeline (Lesson 8).  The stages fetch, decode,
% schedule, execute and write-back are connected by FIFO queues
% (FE, DE and EX queue); the scheduler is drawn as a pool from which
% ready instructions are picked, not as a FIFO.  Execute talks to the
% memory model; completion wakes up dependent instructions.
% Redraws the outline figure m8l1_queue.
% Usage: % Queue-based model of a pipeline (Lesson 8).  The stages fetch, decode,
% schedule, execute and write-back are connected by FIFO queues
% (FE, DE and EX queue); the scheduler is drawn as a pool from which
% ready instructions are picked, not as a FIFO.  Execute talks to the
% memory model; completion wakes up dependent instructions.
% Redraws the outline figure m8l1_queue.
% Usage: \input{figures/l08-queue-model}
% Width: body (120 mm); the drawing is about 116 mm wide.
\begin{tikzpicture}[x=1mm, y=1mm, font=\scriptsize\sffamily,
    >={Stealth[length=1.5mm]},
    st/.style={draw, fill=ln-accent!15, minimum width=14mm,
               minimum height=8mm, align=center, inner sep=1pt},
    q/.style={draw, fill=ln-box, minimum width=12mm, minimum height=4mm,
              inner sep=0pt}]
  % \fifo{x}{y}{name}{label}: a queue drawn as four slots
  \newcommand{\fifo}[4]{%
    \foreach \i in {0,...,3} {
      \draw[fill=ln-box] (#1+\i*3,#2-2) rectangle ++(3,4);
    }
    \coordinate (#3-in) at (#1,#2);
    \coordinate (#3-out) at (#1+12,#2);
    \node[above, inner sep=1pt] at (#1+6,#2+2) {#4};}
  \node[st] (fe) at (7,20) {\iflangja{フェッチ}{fetch}};
  \fifo{18}{20}{feq}{FE queue}
  \node[st] (de) at (41,20) {\iflangja{デコード}{decode}};
  \fifo{52}{20}{deq}{DE queue}
  % scheduler: a pool, not a FIFO
  \node[draw, fill=ln-caution!12, rounded corners=2mm, minimum width=18mm,
        minimum height=10mm, align=center, inner sep=1pt] (sch) at (79,20)
    {\iflangja{スケジューラ}{scheduler}\\[-0.3mm]%
     \iflangja{（FIFO ではない）}{(not a FIFO)}};
  \node[st] (ex) at (103,20) {\iflangja{実行}{execute}};
  \node[st, fill=ln-tip!15] (mem) at (103,4)
    {\iflangja{キャッシュ・メモリ}{caches, memory}};
  \fifo{70}{-2}{exq}{}
  \node[below, inner sep=1pt] at (76,-4) {EX queue};
  \node[st] (wb) at (49,-2) {\iflangja{ライトバック}{write-back}};
  \draw[->] (fe) -- (feq-in);
  \draw[->] (feq-out) -- (de);
  \draw[->] (de) -- (deq-in);
  \draw[->] (deq-out) -- (sch);
  \draw[->] (sch) -- (ex);
  \draw[<->] (ex) -- (mem);
  \draw[->] (ex.east) -- ++(4,0) |- (exq-out);
  \draw[->] (exq-in) -- (wb);
  % wake-up of dependants
  \draw[->, dashed, ln-caution] (wb.north) -- ++(0,5) -| (sch.south)
    node[pos=0.25, above, text=ln-caution, inner sep=1.5pt]
    {\iflangja{完了で依存命令を起こす}{completion wakes dependants}};
  % legend: depth and width
  \draw[|-|, ln-accent] (18,27) -- (30,27)
    node[midway, above, text=ln-accent, inner sep=1.5pt]
    {\iflangja{深さ = ステージの遅延}{depth = stage latency}};
  \draw[->, thick, ln-accent] (45,27) -- (53,27)
    node[right, text=ln-accent, inner sep=1.5pt]
    {\iflangja{幅 = 1サイクルに移動できる命令数}{width = instructions moved per cycle}};
\end{tikzpicture}

% Width: body (120 mm); the drawing is about 116 mm wide.
\begin{tikzpicture}[x=1mm, y=1mm, font=\scriptsize\sffamily,
    >={Stealth[length=1.5mm]},
    st/.style={draw, fill=ln-accent!15, minimum width=14mm,
               minimum height=8mm, align=center, inner sep=1pt},
    q/.style={draw, fill=ln-box, minimum width=12mm, minimum height=4mm,
              inner sep=0pt}]
  % \fifo{x}{y}{name}{label}: a queue drawn as four slots
  \newcommand{\fifo}[4]{%
    \foreach \i in {0,...,3} {
      \draw[fill=ln-box] (#1+\i*3,#2-2) rectangle ++(3,4);
    }
    \coordinate (#3-in) at (#1,#2);
    \coordinate (#3-out) at (#1+12,#2);
    \node[above, inner sep=1pt] at (#1+6,#2+2) {#4};}
  \node[st] (fe) at (7,20) {\iflangja{フェッチ}{fetch}};
  \fifo{18}{20}{feq}{FE queue}
  \node[st] (de) at (41,20) {\iflangja{デコード}{decode}};
  \fifo{52}{20}{deq}{DE queue}
  % scheduler: a pool, not a FIFO
  \node[draw, fill=ln-caution!12, rounded corners=2mm, minimum width=18mm,
        minimum height=10mm, align=center, inner sep=1pt] (sch) at (79,20)
    {\iflangja{スケジューラ}{scheduler}\\[-0.3mm]%
     \iflangja{（FIFO ではない）}{(not a FIFO)}};
  \node[st] (ex) at (103,20) {\iflangja{実行}{execute}};
  \node[st, fill=ln-tip!15] (mem) at (103,4)
    {\iflangja{キャッシュ・メモリ}{caches, memory}};
  \fifo{70}{-2}{exq}{}
  \node[below, inner sep=1pt] at (76,-4) {EX queue};
  \node[st] (wb) at (49,-2) {\iflangja{ライトバック}{write-back}};
  \draw[->] (fe) -- (feq-in);
  \draw[->] (feq-out) -- (de);
  \draw[->] (de) -- (deq-in);
  \draw[->] (deq-out) -- (sch);
  \draw[->] (sch) -- (ex);
  \draw[<->] (ex) -- (mem);
  \draw[->] (ex.east) -- ++(4,0) |- (exq-out);
  \draw[->] (exq-in) -- (wb);
  % wake-up of dependants
  \draw[->, dashed, ln-caution] (wb.north) -- ++(0,5) -| (sch.south)
    node[pos=0.25, above, text=ln-caution, inner sep=1.5pt]
    {\iflangja{完了で依存命令を起こす}{completion wakes dependants}};
  % legend: depth and width
  \draw[|-|, ln-accent] (18,27) -- (30,27)
    node[midway, above, text=ln-accent, inner sep=1.5pt]
    {\iflangja{深さ = ステージの遅延}{depth = stage latency}};
  \draw[->, thick, ln-accent] (45,27) -- (53,27)
    node[right, text=ln-accent, inner sep=1.5pt]
    {\iflangja{幅 = 1サイクルに移動できる命令数}{width = instructions moved per cycle}};
\end{tikzpicture}

% Width: body (120 mm); the drawing is about 116 mm wide.
\begin{tikzpicture}[x=1mm, y=1mm, font=\scriptsize\sffamily,
    >={Stealth[length=1.5mm]},
    st/.style={draw, fill=ln-accent!15, minimum width=14mm,
               minimum height=8mm, align=center, inner sep=1pt},
    q/.style={draw, fill=ln-box, minimum width=12mm, minimum height=4mm,
              inner sep=0pt}]
  % \fifo{x}{y}{name}{label}: a queue drawn as four slots
  \newcommand{\fifo}[4]{%
    \foreach \i in {0,...,3} {
      \draw[fill=ln-box] (#1+\i*3,#2-2) rectangle ++(3,4);
    }
    \coordinate (#3-in) at (#1,#2);
    \coordinate (#3-out) at (#1+12,#2);
    \node[above, inner sep=1pt] at (#1+6,#2+2) {#4};}
  \node[st] (fe) at (7,20) {\iflangja{フェッチ}{fetch}};
  \fifo{18}{20}{feq}{FE queue}
  \node[st] (de) at (41,20) {\iflangja{デコード}{decode}};
  \fifo{52}{20}{deq}{DE queue}
  % scheduler: a pool, not a FIFO
  \node[draw, fill=ln-caution!12, rounded corners=2mm, minimum width=18mm,
        minimum height=10mm, align=center, inner sep=1pt] (sch) at (79,20)
    {\iflangja{スケジューラ}{scheduler}\\[-0.3mm]%
     \iflangja{（FIFO ではない）}{(not a FIFO)}};
  \node[st] (ex) at (103,20) {\iflangja{実行}{execute}};
  \node[st, fill=ln-tip!15] (mem) at (103,4)
    {\iflangja{キャッシュ・メモリ}{caches, memory}};
  \fifo{70}{-2}{exq}{}
  \node[below, inner sep=1pt] at (76,-4) {EX queue};
  \node[st] (wb) at (49,-2) {\iflangja{ライトバック}{write-back}};
  \draw[->] (fe) -- (feq-in);
  \draw[->] (feq-out) -- (de);
  \draw[->] (de) -- (deq-in);
  \draw[->] (deq-out) -- (sch);
  \draw[->] (sch) -- (ex);
  \draw[<->] (ex) -- (mem);
  \draw[->] (ex.east) -- ++(4,0) |- (exq-out);
  \draw[->] (exq-in) -- (wb);
  % wake-up of dependants
  \draw[->, dashed, ln-caution] (wb.north) -- ++(0,5) -| (sch.south)
    node[pos=0.25, above, text=ln-caution, inner sep=1.5pt]
    {\iflangja{完了で依存命令を起こす}{completion wakes dependants}};
  % legend: depth and width
  \draw[|-|, ln-accent] (18,27) -- (30,27)
    node[midway, above, text=ln-accent, inner sep=1.5pt]
    {\iflangja{深さ = ステージの遅延}{depth = stage latency}};
  \draw[->, thick, ln-accent] (45,27) -- (53,27)
    node[right, text=ln-accent, inner sep=1.5pt]
    {\iflangja{幅 = 1サイクルに移動できる命令数}{width = instructions moved per cycle}};
\end{tikzpicture}

  \caption{パイプラインのキューベースモデル．フェッチ，デコード，
    スケジュール，実行，ライトバックは FIFO キューで結ばれ，スケジューラは
    準備のできた命令を選び出すプールである．実行ステージはメモリの
    レイテンシをキャッシュとメモリのモデルから得て，完了した命令は
    それに依存する命令を起こす．}
  \label{fig:l08-queue-model}
\end{figure}

モデルにタイミングを与えるものは3つある．第1は依存関係である．命令が
完了すると，その結果はそれを待つ命令に通知されなければならず，待って
いた命令は準備完了になる．これは第4章のスコアボードがハードウェアで
行うことと同じである．第2に，メモリ命令のレイテンシはキャッシュと
メモリのモデルから得られ，そのモデルが各レベルでのヒットかミスかを
決める．第3に，その他の命令のレイテンシはパラメータであり，たとえば
整数加算に\num{1}~サイクル，乗算に\num{3}~サイクルといった値を，
文献，CACTI のような回路モデル，あるいは設計のレジスタ転送レベル
（RTL）記述から得る．シミュレータ自身は回路をシミュレーションしない．\margin{CACTI
は，キャッシュやレジスタファイルのような SRAM アレイのアクセス時間，
面積，エネルギーを，その容量と構成から求める広く使われたモデルである．}
したがって，実行に依存しない実行レイテンシは参照表である．

\subsection{シミュレータの構造}

サイクルレベルシミュレータは，1回の反復でクロックを1サイクル進め，
すべての構成要素に1サイクル分の仕事をさせるループである．構成要素は
階層をなす．シミュレータはコア，メモリシステム，ネットワークオンチップを
進め，各コアはそのパイプラインステージを進める．\margin{Georgia Tech の
CPU・GPU シミュレータである MacSim \cite{gthparch2026macsim} はこのように
構成されている．リストは我々自身のスケッチであり，そのコードではない．}

\begin{lstlisting}[language=C++, numbers=none]
// 最上位：1回の反復 = シミュレーション上の1サイクル
while (!sim.all_done()) {
  ++sim.cycle;
  for (Core &c : sim.cores) c.tick();
  sim.memory.tick();
  sim.noc.tick();
}

// 1つのコア：ステージをパイプラインの逆順に呼ぶ
void Core::tick() {
  writeback(); execute(); schedule();
  decode(); fetch();
}
\end{lstlisting}

ステージは最後のものから最初のものへ向かって呼ばれる．
\texttt{execute()} が動くとき，それは \texttt{schedule()} が\emph{前の}
サイクルに出力した命令を，\texttt{schedule()} がそれをこのサイクルの
出力で上書きする前に消費する．パイプラインの順に呼ぶと，1つの命令が
1サイクルのうちに5つのステージすべてを通過してしまう．\tip{逆順は，
ステージ間のパイプラインレジスタのソフトウェア上の等価物である．
各ステージは，前のステージが1サイクル前に書いたものを読む．}

シミュレーションされる各命令は小さなレコード（\emph{op}）で表され，
それがタイミングの状態とともに命令をステージ間で運ぶ．鍵となる
フィールドは命令が完了するサイクルである．
\begin{equation}
  \textit{done\_cycle} = \textit{schedule\_cycle} + \textit{latency}.
  \label{eq:l08-done-cycle}
\end{equation}
レイテンシが固定の命令であれば，スケジューラは発行の時点でこれを
計算できる．ロードのレイテンシは固定ではない．アクセスが L1
キャッシュでヒットするか，L2 キャッシュでヒットするか，どちらでも
ミスするかに依存するからである．それはキャッシュ階層でアクセスを
シミュレーションした後に初めて判明し，そのとき発行サイクルに加えられる．

\begin{example}
  \label{ex:l08-done-cycle}
  \texttt{add} に\num{1}~サイクル，\texttt{mul} に\num{3}，L1 でヒット，
  L2 でヒット，両方でミスするロードにそれぞれ\num{4}，\num{30}，
  \num{200}~サイクル（我々自身の数値）というレイテンシのもとで，
  シミュレータは4つの命令を\cref{tab:l08-done-cycle}のように
  スケジュールする．\texttt{add} はロードに依存するので，ロードの
  \textit{done\_cycle} になって初めて準備完了になる．
\end{example}

\begin{table}[tb]
  \centering
  \caption{シミュレータが4つの命令について計算する完了時刻（我々自身の
    レイテンシ）．ロードのレイテンシは，キャッシュモデルがどこでヒット
    するかを決めた後に初めて判明する．}
  \label{tab:l08-done-cycle}
  \begin{tabular}{@{}l c c c >{\raggedright\arraybackslash}p{30mm}@{}}% ja: narrower last column, wider headers
    \toprule
    命令 & スケジュール & レイテンシ & 完了 & レイテンシの出所 \\
    \midrule
    \texttt{mul r1, r2, r3} & 10 & 3 & 13 & 表 \\
    \texttt{ld r4, [r5]} & 11 & 30 & 41 & キャッシュモデル：L1 ミス，L2 ヒット \\
    \texttt{mul r6, r1, r1} & 13 & 3 & 16 & 表；\texttt{r1} を待つ \\
    \texttt{add r7, r4, 1} & 41 & 1 & 42 & 表；\texttt{r4} を待つ \\
    \bottomrule
  \end{tabular}
\end{table}

スケジューラはモデルの中で最も詳細な部分である．各サイクルに，それは
準備のできた命令を，その命令が必要とする資源の制約のもとで選ぶ．空いた
浮動小数点ユニット，ロードキューまたはストアキューのエントリ，空いた
実行ポートやキャッシュポートである．あるユニットを数サイクル占有する命令は
その間ずっと資源を保持し，スケジューラはその間に2つ目の命令をそこへ
発行してはならない．

\subsection{GPU のモデル化}

サイクルレベルの GPU シミュレータも同じ構造をとるが，単位が1つ変わる．
スレッドではなく\emph{ワープ}をモデル化する．SIMT ハードウェアは
1つのワープ命令をそのすべてのスレッドについてまとめてフェッチ，
デコード，スケジュール，実行するからである（第4章）．スケジューラは
各ワープの先頭にある命令を検討する．1つのワープの命令は順に発行され，
異なるワープの命令は任意の順に発行される．\margin{GPU シミュレータが
1サイクルにモデル化すべきものは CPU シミュレータよりはるかに多い．
各 SM は最大64の常駐ワープを保持でき（CPG Table~21
\cite{nvidia2024cudaguide}），GPU には数十の SM がある．}CPU に
対応物のないものが3つある．ダイバージェントなワープ
（\cref{sec:l08-divergence-memory}），ワープ，スレッドブロック，
カーネルの階層，そして第7章のワープスケジューリングポリシーである．

この階層は，シミュレーションがいつ終わるかも決める．スレッドブロック
全体が1つの SM に割り当てられ（第3章），シミュレータはすべての
スレッドを最後の命令を実行するまで追跡する．\margin{SM が一度に保持する
ブロック数もモデル化される制限である．SM あたりの常駐ブロックは
Compute Capability 7.0，8.0，9.0 で最大32，7.5 と 8.6 で16である
（CPG Table~21 \cite{nvidia2024cudaguide}）．}ブロックはそのすべての
スレッドが終了したときに終了し，その SM の資源を次のブロックのために
解放する．カーネルはそのすべてのブロックが終了したときに終了し，
アプリケーションはそのすべてのカーネルが終了したときに終了する．

\begin{keypoint}
  サイクルレベルシミュレータはグローバルなクロック（クロックドメインごとに
  1つ）を進め，各サイクルに各ユニットをモデル化する．命令の流れは，
  プログラムを実行するか，トレースを再生して得る．パイプラインは，深さが
  ステージのレイテンシ，転送速度がパイプライン幅である FIFO キューの
  集合に，資源制約のもとで準備のできた命令を選ぶスケジューラを加えた
  ものである．GPU シミュレータはスレッドではなくワープをモデル化する．
\end{keypoint}

\section{GPU シミュレータにおけるダイバージェンスとメモリ}
\label{sec:l08-divergence-memory}

GPU シミュレータのモデル化の単位はワープである．命令キャッシュへの
アクセスとプログラムカウンタは，スレッドではなくワープに属する．\source{M8，
サイクルレベルシミュレーション（パート~2）；BPG
§10.2.1 Coalesced Access to Global Memory
\cite{nvidia2026cudabestpractices}．}ワープのすべてのスレッドがすべての
命令に参加するわけではないので，シミュレータは第4章のアクティブマスクも
保持し，結果の計算と，命令が実際に使う資源の計数の両方に用いる．本節では，
これまでの章で扱った GPU 固有の2つの効果，ダイバージェンスと
コアレッシングをシミュレータがどう扱うかを示し，続いてオープンソースの
GPU シミュレータを列挙する．

\subsection{ダイバージェント分岐}

シミュレータがダイバージェント分岐（第6章）をどう扱うかは，命令を
どう得るかによる．
\begin{description}
  \item[実行駆動．]  シミュレータは各スレッドの分岐結果を計算し，
    ハードウェアとまったく同じように SIMT スタックを自ら管理する．
    スタックがモデルの一部なので，それを置き換えることで別の機構
    （動的ワープ形成，別の再収束規則）をモデル化できる．
  \item[トレース駆動．]  トレースは，ワープが実行したすべての経路を，
    各基本ブロックのアクティブマスクとともに記録しなければならない．
    たとえば第6章の分岐なら BB1 \texttt{1111}，BB2 \texttt{1100}，
    BB3 \texttt{0011}，BB4 \texttt{1111} である．シミュレータはこれらの
    ブロックを記録された順に再生する．
\end{description}
\caution{トレースは，それを収集したマシンのダイバージェンスの振る舞いを
記録する．そのようなトレースから別のダイバージェンス処理機構を
シミュレーションするのは難しい．トレースには分岐条件も，新しい機構が
選ぶはずの順序も含まれていないからである．}%
したがってトレース駆動シミュレータは，ダイバージェンスの処理そのものを
調べるのには向かない．ダイバージェンスがタイミングを形づくる効果の
1つにすぎない場合には十分である．

\subsection{コアレッシング}

下流のすべて，すなわちキャッシュアクセス，ミス，DRAM トラフィック，
帯域を左右するのは，スレッドの数ではなく，ワープ命令が生成する
メモリトランザクションの数である．したがってメモリコアレッシング
（第4章）のモデル化は GPU シミュレータの最も重要な部分の1つである．
要求は，それが最初に到達するキャッシュの粒度，通常は1キャッシュ
ラインの単位で併合される．本節の小さな構成，すなわち \SI{64}{\byte}
のラインとワープあたり8スレッドでは，スレッドが連続する
\SI{4}{\byte} のワードを読むロード \texttt{R1} は1つの \SI{64}{\byte}
の要求に併合されるが，スレッドが \SI{128}{\byte} ずつ離れたワードを
読むロード \texttt{R2} は，スレッドごとに1つ，それぞれ1ライン全体に
対する8つの要求を生む．

トレースはワープ命令のメモリアクセスを2通りの方法で表せる．すべての
スレッドのアドレス（\texttt{0x00}，\texttt{0x04}，\texttt{0x08}，\dots）
を列挙してコアレッシングをシミュレータに任せるか，コアレッシングの
結果，すなわちトランザクションごとに1つのベースアドレスとサイズを
格納するかである．後者ははるかに小さいが，トレースを収集したマシンの
トランザクションサイズを固定してしまう．\caution{コアレッシング済みの
トレースは1つのマイクロアーキテクチャに縛られる．別のラインサイズや
トランザクションサイズのキャッシュを調べるには，新しいトレースを収集
しなければならない．}

\subsection{キャッシュ階層とセクタキャッシュ}

コアレッシングの後，要求はキャッシュ階層に入り，それは CPU の場合と
ほぼ同じようにモデル化される．アドレスは TLB で変換され（ミス時には
ページテーブルウォーク，第7章），続いてモデル化するマイクロ
アーキテクチャに応じて L1 キャッシュと L2 キャッシュにアクセスする．
異なるのは，その前段にあるコアレッシングの重要性と，キャッシュ自体の
1つの特徴である．

\begin{definition}[セクタキャッシュ]
  \label{def:l08-sectored-cache}
  \term[せくたきゃっしゅ]{セクタキャッシュ}[sectored cache]は，各
  キャッシュブロックを複数のセクタに分割する．各セクタは固有の有効
  ビットを持ち，すべてがブロックの1つのタグを共有する．
  \term[せくた]{セクタ}[sector]はミス時にメモリから埋められる単位で
  あり，キャッシュはブロック全体ではなく，要求されたデータを含む
  セクタだけを取り込む（\cref{fig:l08-sector-cache}）．
\end{definition}

\begin{figure}[tb]
  \centering
  % Coalescing and sectored caches in a simulator (Lesson 8).  Top: one
% 64-byte cache block of a sectored cache, four 16-byte sectors sharing
% one tag, each with its own valid bit.  Bottom: the uncoalesced load R2
% (eight threads, addresses 128 B apart) misses in eight blocks; a
% conventional cache fills eight whole blocks (512 B), a sectored cache
% only the eight sectors that hold the requested words (128 B).
% Redraws the outline figures m8l2_memory_coalescing and
% m8l2_sector_cache (toy sizes of the lecture: 64 B block, 16 B sector;
% sector address typo corrected).
% Usage: % Coalescing and sectored caches in a simulator (Lesson 8).  Top: one
% 64-byte cache block of a sectored cache, four 16-byte sectors sharing
% one tag, each with its own valid bit.  Bottom: the uncoalesced load R2
% (eight threads, addresses 128 B apart) misses in eight blocks; a
% conventional cache fills eight whole blocks (512 B), a sectored cache
% only the eight sectors that hold the requested words (128 B).
% Redraws the outline figures m8l2_memory_coalescing and
% m8l2_sector_cache (toy sizes of the lecture: 64 B block, 16 B sector;
% sector address typo corrected).
% Usage: % Coalescing and sectored caches in a simulator (Lesson 8).  Top: one
% 64-byte cache block of a sectored cache, four 16-byte sectors sharing
% one tag, each with its own valid bit.  Bottom: the uncoalesced load R2
% (eight threads, addresses 128 B apart) misses in eight blocks; a
% conventional cache fills eight whole blocks (512 B), a sectored cache
% only the eight sectors that hold the requested words (128 B).
% Redraws the outline figures m8l2_memory_coalescing and
% m8l2_sector_cache (toy sizes of the lecture: 64 B block, 16 B sector;
% sector address typo corrected).
% Usage: \input{figures/l08-sector-cache}
% Width: body (120 mm); the drawing is about 114 mm wide.
\begin{tikzpicture}[x=1mm, y=1mm, font=\scriptsize\sffamily]
  % one block: tag + 4 sectors with valid bits
  \node[anchor=west, font=\scriptsize\sffamily\bfseries] at (0,51)
    {\iflangja{(a) セクタキャッシュのブロック（64\,B）}{(a) one block of a sectored cache (64\,B)}};
  \draw[fill=ln-accent!15] (0,40) rectangle (16,46) node[midway] {\iflangja{タグ}{tag}};
  \foreach \i in {0,...,3} {
    \draw[fill=ln-box] (17+\i*24,40) rectangle ++(4,6) node[midway] {V};
    \draw[fill=ln-tip!15] (21+\i*24,40) rectangle ++(19,6)
      node[midway] {\iflangja{セクタ}{sector} \i: 16\,B};
  }
  \node[anchor=west, text=ln-muted] at (0,36.5)
    {\iflangja{4つのセクタが1つのタグを共有し，セクタごとに有効ビットを持つ}{four sectors share one tag; each sector has its own valid bit}};
  % R2: eight misses
  \node[anchor=west, font=\scriptsize\sffamily\bfseries] at (0,31)
    {\iflangja{(b) R2 のミス（アドレス 0, 128, \dots, 896）}{(b) misses of R2 (addresses 0, 128, \dots, 896)}};
  \node[anchor=west] at (0,21) {\iflangja{従来型}{conventional}};
  \node[anchor=west] at (0,13) {\iflangja{セクタ型}{sectored}};
  \foreach \k in {0,...,7} {
    \pgfmathsetmacro{\x}{20+\k*11.5}
    \node[text=ln-muted] at (\x+5,26) {\pgfmathparse{int(\k*128)}\pgfmathresult};
    \foreach \s in {0,...,3} {
      \draw[fill=ln-caution!25] (\x+\s*2.5,18) rectangle ++(2.5,6);
      \ifnum\s=0
        \draw[fill=ln-caution!25] (\x+\s*2.5,10) rectangle ++(2.5,6);
      \else
        \draw[fill=white] (\x+\s*2.5,10) rectangle ++(2.5,6);
      \fi
    }
  }
  \node[anchor=west] at (20,5.5)
    {\iflangja{従来型: $8 \times 64 = 512$\,B を転送}{conventional: $8 \times 64 = 512$\,B fetched}};
  \node[anchor=west] at (68,5.5)
    {\iflangja{セクタ型: $8 \times 16 = 128$\,B を転送}{sectored: $8 \times 16 = 128$\,B fetched}};
  \node[anchor=west, text=ln-muted] at (20,1)
    {\iflangja{使われるのは各ブロックの 4\,B（合計 32\,B）}{each block supplies one 4\,B word (32\,B in total)}};
\end{tikzpicture}

% Width: body (120 mm); the drawing is about 114 mm wide.
\begin{tikzpicture}[x=1mm, y=1mm, font=\scriptsize\sffamily]
  % one block: tag + 4 sectors with valid bits
  \node[anchor=west, font=\scriptsize\sffamily\bfseries] at (0,51)
    {\iflangja{(a) セクタキャッシュのブロック（64\,B）}{(a) one block of a sectored cache (64\,B)}};
  \draw[fill=ln-accent!15] (0,40) rectangle (16,46) node[midway] {\iflangja{タグ}{tag}};
  \foreach \i in {0,...,3} {
    \draw[fill=ln-box] (17+\i*24,40) rectangle ++(4,6) node[midway] {V};
    \draw[fill=ln-tip!15] (21+\i*24,40) rectangle ++(19,6)
      node[midway] {\iflangja{セクタ}{sector} \i: 16\,B};
  }
  \node[anchor=west, text=ln-muted] at (0,36.5)
    {\iflangja{4つのセクタが1つのタグを共有し，セクタごとに有効ビットを持つ}{four sectors share one tag; each sector has its own valid bit}};
  % R2: eight misses
  \node[anchor=west, font=\scriptsize\sffamily\bfseries] at (0,31)
    {\iflangja{(b) R2 のミス（アドレス 0, 128, \dots, 896）}{(b) misses of R2 (addresses 0, 128, \dots, 896)}};
  \node[anchor=west] at (0,21) {\iflangja{従来型}{conventional}};
  \node[anchor=west] at (0,13) {\iflangja{セクタ型}{sectored}};
  \foreach \k in {0,...,7} {
    \pgfmathsetmacro{\x}{20+\k*11.5}
    \node[text=ln-muted] at (\x+5,26) {\pgfmathparse{int(\k*128)}\pgfmathresult};
    \foreach \s in {0,...,3} {
      \draw[fill=ln-caution!25] (\x+\s*2.5,18) rectangle ++(2.5,6);
      \ifnum\s=0
        \draw[fill=ln-caution!25] (\x+\s*2.5,10) rectangle ++(2.5,6);
      \else
        \draw[fill=white] (\x+\s*2.5,10) rectangle ++(2.5,6);
      \fi
    }
  }
  \node[anchor=west] at (20,5.5)
    {\iflangja{従来型: $8 \times 64 = 512$\,B を転送}{conventional: $8 \times 64 = 512$\,B fetched}};
  \node[anchor=west] at (68,5.5)
    {\iflangja{セクタ型: $8 \times 16 = 128$\,B を転送}{sectored: $8 \times 16 = 128$\,B fetched}};
  \node[anchor=west, text=ln-muted] at (20,1)
    {\iflangja{使われるのは各ブロックの 4\,B（合計 32\,B）}{each block supplies one 4\,B word (32\,B in total)}};
\end{tikzpicture}

% Width: body (120 mm); the drawing is about 114 mm wide.
\begin{tikzpicture}[x=1mm, y=1mm, font=\scriptsize\sffamily]
  % one block: tag + 4 sectors with valid bits
  \node[anchor=west, font=\scriptsize\sffamily\bfseries] at (0,51)
    {\iflangja{(a) セクタキャッシュのブロック（64\,B）}{(a) one block of a sectored cache (64\,B)}};
  \draw[fill=ln-accent!15] (0,40) rectangle (16,46) node[midway] {\iflangja{タグ}{tag}};
  \foreach \i in {0,...,3} {
    \draw[fill=ln-box] (17+\i*24,40) rectangle ++(4,6) node[midway] {V};
    \draw[fill=ln-tip!15] (21+\i*24,40) rectangle ++(19,6)
      node[midway] {\iflangja{セクタ}{sector} \i: 16\,B};
  }
  \node[anchor=west, text=ln-muted] at (0,36.5)
    {\iflangja{4つのセクタが1つのタグを共有し，セクタごとに有効ビットを持つ}{four sectors share one tag; each sector has its own valid bit}};
  % R2: eight misses
  \node[anchor=west, font=\scriptsize\sffamily\bfseries] at (0,31)
    {\iflangja{(b) R2 のミス（アドレス 0, 128, \dots, 896）}{(b) misses of R2 (addresses 0, 128, \dots, 896)}};
  \node[anchor=west] at (0,21) {\iflangja{従来型}{conventional}};
  \node[anchor=west] at (0,13) {\iflangja{セクタ型}{sectored}};
  \foreach \k in {0,...,7} {
    \pgfmathsetmacro{\x}{20+\k*11.5}
    \node[text=ln-muted] at (\x+5,26) {\pgfmathparse{int(\k*128)}\pgfmathresult};
    \foreach \s in {0,...,3} {
      \draw[fill=ln-caution!25] (\x+\s*2.5,18) rectangle ++(2.5,6);
      \ifnum\s=0
        \draw[fill=ln-caution!25] (\x+\s*2.5,10) rectangle ++(2.5,6);
      \else
        \draw[fill=white] (\x+\s*2.5,10) rectangle ++(2.5,6);
      \fi
    }
  }
  \node[anchor=west] at (20,5.5)
    {\iflangja{従来型: $8 \times 64 = 512$\,B を転送}{conventional: $8 \times 64 = 512$\,B fetched}};
  \node[anchor=west] at (68,5.5)
    {\iflangja{セクタ型: $8 \times 16 = 128$\,B を転送}{sectored: $8 \times 16 = 128$\,B fetched}};
  \node[anchor=west, text=ln-muted] at (20,1)
    {\iflangja{使われるのは各ブロックの 4\,B（合計 32\,B）}{each block supplies one 4\,B word (32\,B in total)}};
\end{tikzpicture}

  \caption{セクタキャッシュ（小さな例：\SI{64}{\byte} のブロック，
    \SI{16}{\byte} のセクタ）．(a) セクタは1つのタグを共有する．
    (b) コアレッシングされないロード \texttt{R2} は8つのブロックで
    ミスする．従来型のキャッシュは8つのブロック全体を埋め，セクタ
    キャッシュはそれぞれ1セクタだけを埋める．}
  \label{fig:l08-sector-cache}
\end{figure}

コアレッシングされたロード \texttt{R1}（連続する8つの \SI{4}{\byte}
ワード，アドレス 0 から 31 の \SI{32}{\byte}）では，従来型のキャッシュは
\SI{64}{\byte} のブロック全体を埋めるが，セクタキャッシュはデータを
含む \SI{32}{\byte}，すなわちセクタ 0 と~1 だけを埋める．節約は2倍で，
取り込んだバイトはすべて使われる．コアレッシングされないロード
\texttt{R2} では，従来型のキャッシュは $8 \times 4 = \SI{32}{\byte}$
の有用なデータを届けるために $8 \times 64 = \SI{512}{\byte}$ を
取り込み，セクタキャッシュは $8 \times 16 = \SI{128}{\byte}$ を取り込む．
したがってセクタ化の利点は散在したアクセスに対するメモリ帯域の削減で
あり，複数のセクタで1つのタグを共有することで，大きなブロックの
タグ記憶量が保たれる．代償は暗黙のプリフェッチが減ることである．
セクタ単位で埋められるブロックは，空間的局所性によってまもなく使われた
はずの隣接データを持ち込まなくなる．\caution{講義 vs ガイド：
\SI{64}{\byte} のブロック中の \SI{16}{\byte} のセクタは小さな例の
サイズである．NVIDIA GPU は \SI{128}{\byte} のラインに
\SI{32}{\byte} のセクタを用い，Compute Capability 6.0 以降，ワープの
アクセスのコストは\enquote{the number of 32-byte transactions necessary
to service all of the threads of the warp}である（BPG §10.2.1
\cite{nvidia2026cudabestpractices}）．第4章はこれらのセクタを数えている．}

\begin{note}
  セクタキャッシュをシミュレーションすると，シミュレータが数えるものが
  変わる．ブロックが一部のセクタだけ有効な状態で存在しうるので，
  アクセスがタグにヒットしながらセクタではミスすることがある．タグの
  ヒットだけを数えるシミュレータは，ヒット率を過大に，DRAM
  トラフィックを過小に見積もる．
\end{note}

\subsection{オープンソースの GPU シミュレータ}

\cref{tab:l08-simulators}は，研究で広く使われているオープンソースの
GPU シミュレータを示す．それらは，モデル化する ISA，命令の得方，
1つの GPU を超えて何をモデル化するかで異なる．

\begin{table}[tb]
  \centering
  \caption{オープンソースの GPU シミュレータ（本章の分類による）．}
  \label{tab:l08-simulators}
  \setlength{\tabcolsep}{4pt}
  \begin{tabular}{@{}>{\raggedright\arraybackslash}p{24mm}
                  >{\raggedright\arraybackslash}p{25mm}
                  >{\raggedright\arraybackslash}p{22mm}% ja: wider for トレース駆動
                  >{\raggedright\arraybackslash}p{37mm}@{}}
    \toprule
    シミュレータ & ISA & 種類 & 対象 \\
    \midrule
    GPGPU-Sim \cite{bakhoda2009gpgpusim} & NVIDIA PTX，SASS
      & 実行駆動 & 汎用コードを実行する GPU \\
    Accel-Sim \cite{khairy2020accelsim} & NVIDIA PTX，SASS
      & トレース駆動 & GPU とアクセラレータ \\
    MGPUSim \cite{sun2019mgpusim} & AMD GPU & 実行駆動
      & 複数の GPU \\
    MacSim \cite{gthparch2026macsim} & NVIDIA と Intel の GPU
      & トレース駆動 & CPU と GPU のヘテロジニアスシステム \\
    gem5-gpu \cite{power2015gem5gpu} & AMD GPU または NVIDIA PTX
      & 実行駆動 & CPU と GPU のヘテロジニアスシステム \\
    \bottomrule
  \end{tabular}
\end{table}

GPGPU-Sim はこれらの中で最も広く使われている．gem5-gpu はそれを最も
広く使われている CPU シミュレータ gem5 と組み合わせ，CPU と GPU が
ともに実行してメモリシステムを共有できるようにする．MGPUSim は AMD の
GPU をモデル化し，\cref{sec:l08-accelerating}の Photon のサンプリング
手法はその上に構築された．\sidenote{GPGPU-Sim は ISPASS 2009 での CUDA
ワークロードの研究とともに発表された \cite{bakhoda2009gpgpusim}．表は
本章の出典の分類に従っている．詳しく言えば，Accel-Sim はトレース駆動の
SASS モードと並んで GPGPU-Sim の実行駆動の PTX モードを残しており
\cite{khairy2020accelsim}，gem5-gpu（\enquote{Gem5-GPGPU-Sim}とも
記される組み合わせ）は gem5 を GPGPU-Sim の NVIDIA モデルと結合する．
gem5 の AMD GPU モデルは別の開発である．MacSim の文書は，それを x86，
ARM64，PTX，Intel GEN の命令に対して\enquote{either a trace driven or
execution-driven cycle level simulator}として構成できるものと説明して
いる \cite{gthparch2026macsim}．}

\begin{keypoint}
  GPU シミュレータはワープをそのアクティブマスクとともにモデル化する．
  SIMT スタックを自ら管理する（実行駆動）か，記録されたブロックと
  マスクを再生する（トレース駆動．トレースを取ったマシンの
  ダイバージェンス処理に縛られる）．コアレッシングはトランザクション数を
  決め，モデル化すべき最も重要なものの1つである．コアレッシング済みの
  トレースはトランザクションサイズを固定する．セクタキャッシュはミス
  ごとに1セクタを埋める．帯域は減るが，暗黙のプリフェッチも減る．
\end{keypoint}

\section{GPU の解析モデル}
\label{sec:l08-analytical}

サイクルレベルシミュレータは，1つの設計と1つのプログラムについての問いに，
そのプログラムを実行して答える．\source{M8，解析モデル；K12 ch.~2
Performance Principles，§4.1.1 Execution Time Modeling，§4.2.4
Roofline Model，§4.2.6 Other GPGPU Performance Modeling Techniques；
\cite{williams2009roofline,hong2009analytical,huang2014gpumech,karkhanis2004firstorder}．}
解析モデルはそれに式で答える．

\begin{definition}[解析モデル]
  \label{def:l08-analytical-model}
  \term[かいせきもでる]{解析モデル}[analytical model]は，設計と
  プログラムの少数のパラメータから，プログラムを実行するのではなく
  方程式によって性能を予測する．
  \term[いちじきんじもでる]{一次近似モデル}[first-order model]は，
  性能を支配する効果だけを残し，二次的な相互作用を無視する．
\end{definition}

解析モデルは完全な実行を必要とせず，推論できるほど単純で，洞察を与える．
設計が\emph{なぜ}速いのか遅いのか，どのパラメータを変えるべきかを
教えてくれる．その精度はシミュレータより低いので，設計空間を探索して
その最良の点をシミュレーションする前段として，またシミュレーション
結果を説明するために使われる．

\subsection{GPU の一次近似設計}

1つの目標に対して GPU の規模を決める．1秒あたり\num{e12}回の
内積ステップであり，各ステップは
\lstinline[language=C++]!sum += x[i] * y[i]! である．

\begin{definition}[融合積和演算]
  \label{def:l08-fma}
  \term[ゆうごうせきわえんざん]{融合積和演算（FMA）}[fused multiply-add (FMA)]
  は，$a \cdot b + c$ を1回の丸めで1つの命令として計算する．これは
  2回の浮動小数点演算（FLOP）として数える．
\end{definition}

\begin{description}
  \item[演算．]  各ステップは1回の乗算と1回の加算なので，目標は
    \SI{2}{\tera\flop\per\second}，すなわち1秒あたり\num{e12}回の
    FMA である．\SI{1}{\giga\hertz} ではこれに\num{1000}個の FMA
    ユニットが必要であり，\SI{2}{\giga\hertz} では\num{500}個である．
  \item[メモリ．]  各ステップは \texttt{x} の要素を1つ，\texttt{y} の
    要素を1つ読む．この設計演習のように要素あたり1バイトと数えると，
    GPU には \SI{2}{\tera\byte\per\second} のメモリ帯域が必要である．\caution{講義
    vs 実際：\enquote{ステップあたり2バイト}は1バイトの要素を
    前提としている．FP32 の要素ではステップは2~FLOP に対して
    \SI{8}{\byte} を読み，同じ設計には \SI{8}{\tera\byte\per\second}
    が必要になる．これは H100 SXM5 の \SI{3}{\tera\byte\per\second}
    （暫定値，\cite{nvidia2022hopper}）の約2.7倍である．}
  \item[ワープ．]  \num{500}個の FMA ユニットは，各レーンが1サイクルに
    1つの FMA を発行する幅32のワープ約16個（$16 \times 32 = 512$）に
    相当する．GPU は，それぞれが1サイクルに1ワープを発行する16個の
    SM としても，それぞれが2ワープを発行する8個の SM としても構成
    できる．どちらも \SI{2}{\giga\hertz} である．
\end{description}

\subsection{マルチスレッディングの幅と深さ}

SM が保持しなければならないワープの数は，2つのパラメータから決まる．

\begin{definition}[$W_{\text{width}}$ と $W_{\text{depth}}$]
  \label{def:l08-wwidth-wdepth}
  $W_{\text{width}}$ は1サイクルに実行するスレッドの数である（たとえば
  1ワープの32）．$W_{\text{depth}}$ は，1つのワープがストールしている間に
  SM がスケジュールできるワープの数であり，これによってストールが
  隠される．SM は $W_{\text{width}} \times W_{\text{depth}}$ 個の
  常駐スレッドを保持する．
\end{definition}

\begin{figure}[tb]
  \centering
  % W_width and W_depth (Lesson 8).  Left: issue slots of one SM for
% 1, 2 and 4 warps, each warp issuing one compute (C) and one memory (M)
% instruction and then waiting until 8 cycles after its compute
% instruction; grey = idle issue slot.  With 4 warps the stall is fully
% hidden (W_depth = 4).  Right: W_width x W_depth resident threads.
% Redraws the outline figure m8l3_arch2 (our own latency of 8 cycles).
% Usage: % W_width and W_depth (Lesson 8).  Left: issue slots of one SM for
% 1, 2 and 4 warps, each warp issuing one compute (C) and one memory (M)
% instruction and then waiting until 8 cycles after its compute
% instruction; grey = idle issue slot.  With 4 warps the stall is fully
% hidden (W_depth = 4).  Right: W_width x W_depth resident threads.
% Redraws the outline figure m8l3_arch2 (our own latency of 8 cycles).
% Usage: % W_width and W_depth (Lesson 8).  Left: issue slots of one SM for
% 1, 2 and 4 warps, each warp issuing one compute (C) and one memory (M)
% instruction and then waiting until 8 cycles after its compute
% instruction; grey = idle issue slot.  With 4 warps the stall is fully
% hidden (W_depth = 4).  Right: W_width x W_depth resident threads.
% Redraws the outline figure m8l3_arch2 (our own latency of 8 cycles).
% Usage: \input{figures/l08-wdepth}
% Width: body (120 mm); the drawing is about 116 mm wide.
\begin{tikzpicture}[x=4.4mm, y=4.4mm, font=\footnotesize\sffamily]
  % \slots{y}{label}{list}: a row of issue slots; "-" is idle
  \def\idle{-}%
  \newcommand{\slots}[3]{%
    \node[anchor=east, inner sep=1pt] at (-0.2,#1-0.5) {#2};
    \foreach \s [count=\i from 0] in {#3} {
      \ifx\s\idle
        \draw[fill=ln-rule!40, draw=ln-rule] (\i,#1) rectangle ++(1,-1);
      \else
        \draw[fill=ln-box] (\i,#1) rectangle ++(1,-1) node[midway] {\s};
      \fi
    }}
  \slots{0}{\iflangja{1ワープ}{1 warp}}{C1,M1,-,-,-,-,-,-,C1,M1,-,-,-,-,-,-}
  \slots{-1.6}{\iflangja{2ワープ}{2 warps}}{C1,M1,C2,M2,-,-,-,-,C1,M1,C2,M2,-,-,-,-}
  \slots{-3.2}{\iflangja{4ワープ}{4 warps}}{C1,M1,C2,M2,C3,M3,C4,M4,C1,M1,C2,M2,C3,M3,C4,M4}
  \draw[->, ln-muted] (0,-4.8) -- (16.3,-4.8)
    node[below left, inner sep=1pt] {\iflangja{サイクル}{cycle}};
  \draw[ln-caution, thick, ->] (0,0.4) -- (8,0.4)
    node[midway, above, inner sep=1pt, text=ln-caution]
    {\iflangja{ワープ1の周期 = 8サイクル}{one round of warp 1 = 8 cycles}};
  % W_width x W_depth grid
  \begin{scope}[shift={(18.6,-4.2)}, x=1.4mm, y=1.4mm]
    \foreach \r in {0,...,3} {
      \foreach \c in {0,...,7} {
        \draw[fill=ln-accent!15] (\c,\r) circle (0.45);
      }
    }
    \draw[->] (0,4.3) -- (7.6,4.3) node[midway, above, inner sep=1pt]
      {$W_{\text{width}}$};
    \draw[->] (8.4,3.4) -- (8.4,-0.4) node[midway, right, inner sep=1pt]
      {$W_{\text{depth}}$};
    \node[anchor=north, align=center, text=ln-muted] at (4,-0.8)
      {\iflangja{常駐スレッド数\\$= W_{\text{width}} \times W_{\text{depth}}$}%
                {resident threads\\$= W_{\text{width}} \times W_{\text{depth}}$}};
  \end{scope}
\end{tikzpicture}

% Width: body (120 mm); the drawing is about 116 mm wide.
\begin{tikzpicture}[x=4.4mm, y=4.4mm, font=\footnotesize\sffamily]
  % \slots{y}{label}{list}: a row of issue slots; "-" is idle
  \def\idle{-}%
  \newcommand{\slots}[3]{%
    \node[anchor=east, inner sep=1pt] at (-0.2,#1-0.5) {#2};
    \foreach \s [count=\i from 0] in {#3} {
      \ifx\s\idle
        \draw[fill=ln-rule!40, draw=ln-rule] (\i,#1) rectangle ++(1,-1);
      \else
        \draw[fill=ln-box] (\i,#1) rectangle ++(1,-1) node[midway] {\s};
      \fi
    }}
  \slots{0}{\iflangja{1ワープ}{1 warp}}{C1,M1,-,-,-,-,-,-,C1,M1,-,-,-,-,-,-}
  \slots{-1.6}{\iflangja{2ワープ}{2 warps}}{C1,M1,C2,M2,-,-,-,-,C1,M1,C2,M2,-,-,-,-}
  \slots{-3.2}{\iflangja{4ワープ}{4 warps}}{C1,M1,C2,M2,C3,M3,C4,M4,C1,M1,C2,M2,C3,M3,C4,M4}
  \draw[->, ln-muted] (0,-4.8) -- (16.3,-4.8)
    node[below left, inner sep=1pt] {\iflangja{サイクル}{cycle}};
  \draw[ln-caution, thick, ->] (0,0.4) -- (8,0.4)
    node[midway, above, inner sep=1pt, text=ln-caution]
    {\iflangja{ワープ1の周期 = 8サイクル}{one round of warp 1 = 8 cycles}};
  % W_width x W_depth grid
  \begin{scope}[shift={(18.6,-4.2)}, x=1.4mm, y=1.4mm]
    \foreach \r in {0,...,3} {
      \foreach \c in {0,...,7} {
        \draw[fill=ln-accent!15] (\c,\r) circle (0.45);
      }
    }
    \draw[->] (0,4.3) -- (7.6,4.3) node[midway, above, inner sep=1pt]
      {$W_{\text{width}}$};
    \draw[->] (8.4,3.4) -- (8.4,-0.4) node[midway, right, inner sep=1pt]
      {$W_{\text{depth}}$};
    \node[anchor=north, align=center, text=ln-muted] at (4,-0.8)
      {\iflangja{常駐スレッド数\\$= W_{\text{width}} \times W_{\text{depth}}$}%
                {resident threads\\$= W_{\text{width}} \times W_{\text{depth}}$}};
  \end{scope}
\end{tikzpicture}

% Width: body (120 mm); the drawing is about 116 mm wide.
\begin{tikzpicture}[x=4.4mm, y=4.4mm, font=\footnotesize\sffamily]
  % \slots{y}{label}{list}: a row of issue slots; "-" is idle
  \def\idle{-}%
  \newcommand{\slots}[3]{%
    \node[anchor=east, inner sep=1pt] at (-0.2,#1-0.5) {#2};
    \foreach \s [count=\i from 0] in {#3} {
      \ifx\s\idle
        \draw[fill=ln-rule!40, draw=ln-rule] (\i,#1) rectangle ++(1,-1);
      \else
        \draw[fill=ln-box] (\i,#1) rectangle ++(1,-1) node[midway] {\s};
      \fi
    }}
  \slots{0}{\iflangja{1ワープ}{1 warp}}{C1,M1,-,-,-,-,-,-,C1,M1,-,-,-,-,-,-}
  \slots{-1.6}{\iflangja{2ワープ}{2 warps}}{C1,M1,C2,M2,-,-,-,-,C1,M1,C2,M2,-,-,-,-}
  \slots{-3.2}{\iflangja{4ワープ}{4 warps}}{C1,M1,C2,M2,C3,M3,C4,M4,C1,M1,C2,M2,C3,M3,C4,M4}
  \draw[->, ln-muted] (0,-4.8) -- (16.3,-4.8)
    node[below left, inner sep=1pt] {\iflangja{サイクル}{cycle}};
  \draw[ln-caution, thick, ->] (0,0.4) -- (8,0.4)
    node[midway, above, inner sep=1pt, text=ln-caution]
    {\iflangja{ワープ1の周期 = 8サイクル}{one round of warp 1 = 8 cycles}};
  % W_width x W_depth grid
  \begin{scope}[shift={(18.6,-4.2)}, x=1.4mm, y=1.4mm]
    \foreach \r in {0,...,3} {
      \foreach \c in {0,...,7} {
        \draw[fill=ln-accent!15] (\c,\r) circle (0.45);
      }
    }
    \draw[->] (0,4.3) -- (7.6,4.3) node[midway, above, inner sep=1pt]
      {$W_{\text{width}}$};
    \draw[->] (8.4,3.4) -- (8.4,-0.4) node[midway, right, inner sep=1pt]
      {$W_{\text{depth}}$};
    \node[anchor=north, align=center, text=ln-muted] at (4,-0.8)
      {\iflangja{常駐スレッド数\\$= W_{\text{width}} \times W_{\text{depth}}$}%
                {resident threads\\$= W_{\text{width}} \times W_{\text{depth}}$}};
  \end{scope}
\end{tikzpicture}

  \caption{メモリのストールの隠蔽．各ワープは演算命令（C）とメモリ命令
    （M）を1つずつ発行し，演算命令の8~サイクル後になって初めて再び
    発行できる．1ワープでは発行スロットは8サイクル中2サイクル，
    2ワープでは8サイクル中4サイクル埋まり，4ワープでは常に埋まる．
    $W_{\text{depth}} = 4$ である．}
  \label{fig:l08-wdepth}
\end{figure}

\cref{fig:l08-wdepth}は，ワープが演算命令とメモリ命令を1つずつ交互に
実行するプログラムについてこの考え方を示す．1ワープでは，SM は2つの
命令を発行した後，メモリ命令が完了するまでアイドルになる．2ワープでは，
2つ目のワープがアイドルなスロットのうち2つを埋める．4ワープでは
アイドルなスロットがなくなる．一般に，ストールは他のワープの仕事で
覆われなければならない．各ワープが $n$ 個の命令を発行した後，その最初の
命令から $L$ サイクル後まで再び発行できないならば，およそ
\begin{equation}
  W_{\text{depth}} \approx \frac{L}{n}
  \label{eq:l08-wdepth}
\end{equation}
個のワープが必要である．\tip{ワープが2回の待ちの間に発行スロットを
占めるサイクル数を数え，待ち時間をそれで割る．}メモリレイテンシが
\num{200}~サイクルで，ワープあたり演算命令1つとメモリ命令1つなら
$W_{\text{depth}} = 200/(1+1) = 100$ であり，演算命令4つあたり
メモリ命令1つなら $200/(1+4) = 40$ である．メモリアクセスあたりの演算が
多いほど，レイテンシを隠すのに必要なワープは少ない．

どちらのパラメータもハードウェアによって決まる．$W_{\text{width}}$ は
ALU の数（とスケジューラの幅）で制限される．$W_{\text{depth}}$ は
ワープの状態を格納する記憶容量で制限される．$W_{\text{depth}}$ が20
なら，スレッドあたり $r$ 個のレジスタについて $20 \times 32 \times r$
個のレジスタと，ワープごとに1つ，少なくとも20個のプログラム
カウンタが必要である．\margin{実際の制限：SM あたりの常駐ワープは最大64
（Compute Capability 7.0，8.0，9.0），SM あたりの32ビットレジスタは
\num{65536}個なので，32スレッドのワープ64個はそれぞれ最大32個の
レジスタを得る（CPG Table~21 \cite{nvidia2024cudaguide}）．}レイテンシが
数百サイクルのときの式~\eqref{eq:l08-wdepth}は，GPU が SM あたり数十の
ワープを保持する理由と，そのレジスタファイルがこれほど大きい理由
（第6章）を説明する．

\subsection{SM の数とサブコア}

上の2つの設計，すなわち1サイクルに1ワープを発行する16個の SM と2ワープを
発行する8個の SM は，実行帯域が同じであり，同じチップ面積なら合計の
キャッシュ容量も同じである．違うのはその容量の分け方である．少数の
大きな SM はアクセス時間の長い大きなキャッシュを持つが，1つの SM 上の
スレッドブロックはより多くのデータを共有し，それがヒット率を上げうる．
多数の小さな SM は小さく速いキャッシュを持つ．どちらがよいかはブロックが
どれだけデータを共有するかに依存し，このトレードオフは一次近似モデルでは
決着がつかない．より細かいモデルかシミュレーションが必要である．

\begin{definition}[サブコア]
  \label{def:l08-sub-core}
  \term[さぶこあ]{サブコア}[sub-core]は SM をいくつかに分割した区画の
  1つであり，それぞれが固有のプログラムカウンタ，デコーダ，ワープ
  スケジューラ，実行レーンを持ち，SM の命令キャッシュと，クロスバを
  介してその大きなキャッシュと共有メモリを共有する
  （\cref{fig:l08-subcore}）．
\end{definition}

\begin{figure}[tb]
  \centering
  % An SM divided into sub-cores (Lesson 8).  Each sub-core has its own
% PC registers, decoder, warp scheduler and execution lanes; all
% sub-cores share the instruction cache at the top and, through a
% crossbar, the large caches at the bottom.
% Redraws the outline figure m8l3_subcore (four sub-cores instead of
% three, as in current NVIDIA SMs).
% Usage: % An SM divided into sub-cores (Lesson 8).  Each sub-core has its own
% PC registers, decoder, warp scheduler and execution lanes; all
% sub-cores share the instruction cache at the top and, through a
% crossbar, the large caches at the bottom.
% Redraws the outline figure m8l3_subcore (four sub-cores instead of
% three, as in current NVIDIA SMs).
% Usage: % An SM divided into sub-cores (Lesson 8).  Each sub-core has its own
% PC registers, decoder, warp scheduler and execution lanes; all
% sub-cores share the instruction cache at the top and, through a
% crossbar, the large caches at the bottom.
% Redraws the outline figure m8l3_subcore (four sub-cores instead of
% three, as in current NVIDIA SMs).
% Usage: \input{figures/l08-subcore}
% Width: body (120 mm); the drawing is about 108 mm wide.
\begin{tikzpicture}[x=1mm, y=1mm, font=\scriptsize\sffamily,
    u/.style={draw, fill=white, minimum width=21mm, minimum height=4.6mm,
              inner sep=0.5pt}]
  \draw[ln-accent, thick, rounded corners=1.5mm] (0,-2) rectangle (108,50);
  \node[anchor=north west, text=ln-accent, font=\scriptsize\sffamily\bfseries]
    at (1,49) {SM};
  \node[draw, fill=ln-tip!15, minimum width=96mm, minimum height=5mm]
    at (57,43) {\iflangja{命令キャッシュ（共有）}{instruction cache (shared)}};
  \foreach \k in {0,...,3} {
    \pgfmathsetmacro{\x}{8+\k*25}
    \draw[fill=ln-accent!10] (\x,11) rectangle ++(23,28);
    \node[anchor=north] at (\x+11.5,38.6) {\iflangja{サブコア}{sub-core} \k};
    \node[u] at (\x+11.5,31) {\iflangja{PC レジスタ}{PC registers}};
    \node[u] at (\x+11.5,26) {\iflangja{デコーダ}{decoder}};
    \node[u] at (\x+11.5,21) {\iflangja{ワープスケジューラ}{warp scheduler}};
    \node[u, fill=ln-caution!12] at (\x+11.5,15.6) {\iflangja{実行レーン}{lanes (SPs)}};
  }
  \node[draw, fill=ln-box, minimum width=96mm, minimum height=4.5mm]
    at (57,6) {\iflangja{クロスバ}{crossbar}};
  \node[draw, fill=ln-tip!15, minimum width=96mm, minimum height=4.5mm]
    at (57,0.8) {\iflangja{共有メモリと L1 キャッシュ（共有）}{shared memory and L1 cache (shared)}};
\end{tikzpicture}

% Width: body (120 mm); the drawing is about 108 mm wide.
\begin{tikzpicture}[x=1mm, y=1mm, font=\scriptsize\sffamily,
    u/.style={draw, fill=white, minimum width=21mm, minimum height=4.6mm,
              inner sep=0.5pt}]
  \draw[ln-accent, thick, rounded corners=1.5mm] (0,-2) rectangle (108,50);
  \node[anchor=north west, text=ln-accent, font=\scriptsize\sffamily\bfseries]
    at (1,49) {SM};
  \node[draw, fill=ln-tip!15, minimum width=96mm, minimum height=5mm]
    at (57,43) {\iflangja{命令キャッシュ（共有）}{instruction cache (shared)}};
  \foreach \k in {0,...,3} {
    \pgfmathsetmacro{\x}{8+\k*25}
    \draw[fill=ln-accent!10] (\x,11) rectangle ++(23,28);
    \node[anchor=north] at (\x+11.5,38.6) {\iflangja{サブコア}{sub-core} \k};
    \node[u] at (\x+11.5,31) {\iflangja{PC レジスタ}{PC registers}};
    \node[u] at (\x+11.5,26) {\iflangja{デコーダ}{decoder}};
    \node[u] at (\x+11.5,21) {\iflangja{ワープスケジューラ}{warp scheduler}};
    \node[u, fill=ln-caution!12] at (\x+11.5,15.6) {\iflangja{実行レーン}{lanes (SPs)}};
  }
  \node[draw, fill=ln-box, minimum width=96mm, minimum height=4.5mm]
    at (57,6) {\iflangja{クロスバ}{crossbar}};
  \node[draw, fill=ln-tip!15, minimum width=96mm, minimum height=4.5mm]
    at (57,0.8) {\iflangja{共有メモリと L1 キャッシュ（共有）}{shared memory and L1 cache (shared)}};
\end{tikzpicture}

% Width: body (120 mm); the drawing is about 108 mm wide.
\begin{tikzpicture}[x=1mm, y=1mm, font=\scriptsize\sffamily,
    u/.style={draw, fill=white, minimum width=21mm, minimum height=4.6mm,
              inner sep=0.5pt}]
  \draw[ln-accent, thick, rounded corners=1.5mm] (0,-2) rectangle (108,50);
  \node[anchor=north west, text=ln-accent, font=\scriptsize\sffamily\bfseries]
    at (1,49) {SM};
  \node[draw, fill=ln-tip!15, minimum width=96mm, minimum height=5mm]
    at (57,43) {\iflangja{命令キャッシュ（共有）}{instruction cache (shared)}};
  \foreach \k in {0,...,3} {
    \pgfmathsetmacro{\x}{8+\k*25}
    \draw[fill=ln-accent!10] (\x,11) rectangle ++(23,28);
    \node[anchor=north] at (\x+11.5,38.6) {\iflangja{サブコア}{sub-core} \k};
    \node[u] at (\x+11.5,31) {\iflangja{PC レジスタ}{PC registers}};
    \node[u] at (\x+11.5,26) {\iflangja{デコーダ}{decoder}};
    \node[u] at (\x+11.5,21) {\iflangja{ワープスケジューラ}{warp scheduler}};
    \node[u, fill=ln-caution!12] at (\x+11.5,15.6) {\iflangja{実行レーン}{lanes (SPs)}};
  }
  \node[draw, fill=ln-box, minimum width=96mm, minimum height=4.5mm]
    at (57,6) {\iflangja{クロスバ}{crossbar}};
  \node[draw, fill=ln-tip!15, minimum width=96mm, minimum height=4.5mm]
    at (57,0.8) {\iflangja{共有メモリと L1 キャッシュ（共有）}{shared memory and L1 cache (shared)}};
\end{tikzpicture}

  \caption{4つのサブコアに分割された SM．それぞれが固有のフロントエンドと
    レーンを持ち，命令キャッシュとデータキャッシュは共有される．}
  \label{fig:l08-subcore}
\end{figure}

サブコアは2つの設計の中間の方法である．各サブコアは自分のレーン上で
少数のワープをスケジュールし，大きな構造は共有されたままである．
モデルにとっては資源管理が加わる．異なるサブコア上の2つのワープは，
共有キャッシュとクロスバを奪い合うが，発行スロットは奪い合わない．\margin{現在の
NVIDIA の SM は128個の FP32 レーンを4つの区画に分割しており
\cite{nvidia2022hopper}，各サブコアは32レーンに対して発行する．}

\begin{keypoint}
  一次近似設計は目標スループットから出発する．FLOP 数が（与えられた
  クロックでの）FMA ユニットの数を決め，FLOP あたりのバイト数が
  メモリ帯域を決め，隠すべきレイテンシが $W_{\text{depth}} \approx L/n$，
  したがって SM あたりのレジスタとプログラムカウンタを決める．
\end{keypoint}

\subsection{ルーフラインモデル}

一次近似設計は，プログラムがどれだけの演算とどれだけの帯域を必要と
するかを問うた．ルーフラインモデルは逆の問いを立てる．マシンが与えられた
とき，その2つのうちどちらがプログラムを制限するか，である．

\begin{definition}[演算強度]
  \label{def:l08-arithmetic-intensity}
  プログラムまたはカーネルの\term[えんざんきょうど]{演算強度}[arithmetic intensity]
  は，それがアクセスするメモリ1バイトあたりに行う浮動小数点演算の数で
  あり，単位は FLOP/B である．
\end{definition}

\begin{definition}[ルーフラインモデル]
  \label{def:l08-roofline}
  \term[るーふらいんもでる]{ルーフラインモデル}[Roofline model]は，
  ピーク演算性能 $P_{\text{peak}}$ とピークメモリ帯域
  $B_{\text{peak}}$ を持つマシン上で，演算強度 $I$ のカーネルが達成
  できる性能を
  \begin{equation}
    P(I) = \min\bigl(P_{\text{peak}},\; B_{\text{peak}} \cdot I\bigr)
    \label{eq:l08-roofline}
  \end{equation}
  で抑える．両対数軸では，この上界は傾き1の斜めの線とそれに続く
  平らな屋根になる（\cref{fig:l08-roofline}）．2本の線が交わる演算強度
  $I^{*} = P_{\text{peak}}/B_{\text{peak}}$ を
  \term[りっじぽいんと]{リッジポイント}[ridge point]と呼び，その左の
  カーネルは\term[めもりりっそく]{メモリ律速}[memory-bound]であって，
  性能は帯域で制限される．その右のカーネルは
  \term[えんざんりっそく]{演算律速}[compute-bound]であり，ピークの
  FLOP レートで制限される．
\end{definition}

\begin{figure}[tb]
  \centering
  % The Roofline model (Lesson 8), log-log axes, for our example machine
% with a peak of 16 TFLOP/s and 1 TB/s of memory bandwidth (ridge point
% 16 FLOP/B).  The slanted roof is bandwidth x intensity (memory-bound
% region), the flat roof the peak (compute-bound region).  Kernel A has
% an arithmetic intensity of 0.5 FLOP/B, kernel B of 2 FLOP/B: raising
% the intensity moves a kernel right, toward the ridge point.
% Redraws the outline figure m8l3_roffline [sic] with our own numbers.
% Usage: % The Roofline model (Lesson 8), log-log axes, for our example machine
% with a peak of 16 TFLOP/s and 1 TB/s of memory bandwidth (ridge point
% 16 FLOP/B).  The slanted roof is bandwidth x intensity (memory-bound
% region), the flat roof the peak (compute-bound region).  Kernel A has
% an arithmetic intensity of 0.5 FLOP/B, kernel B of 2 FLOP/B: raising
% the intensity moves a kernel right, toward the ridge point.
% Redraws the outline figure m8l3_roffline [sic] with our own numbers.
% Usage: % The Roofline model (Lesson 8), log-log axes, for our example machine
% with a peak of 16 TFLOP/s and 1 TB/s of memory bandwidth (ridge point
% 16 FLOP/B).  The slanted roof is bandwidth x intensity (memory-bound
% region), the flat roof the peak (compute-bound region).  Kernel A has
% an arithmetic intensity of 0.5 FLOP/B, kernel B of 2 FLOP/B: raising
% the intensity moves a kernel right, toward the ridge point.
% Redraws the outline figure m8l3_roffline [sic] with our own numbers.
% Usage: \input{figures/l08-roofline}
% Width: body (120 mm); the drawing is about 110 mm wide.
\begin{tikzpicture}[x=8.5mm, y=6mm, font=\footnotesize\sffamily,
    >={Stealth[length=1.5mm]}]
  % regions (x = log2 FLOP/B from -3 to 7, y = log2 TFLOP/s from -3 to 5)
  \fill[ln-accent!8] (-3,-3) rectangle (4,5.4);
  \fill[ln-tip!10] (4,-3) rectangle (7,5.4);
  \node[text=ln-accent, align=center] at (-0.2,4.6)
    {\iflangja{メモリ律速}{memory-bound}};
  \node[text=ln-tip, align=center] at (5.5,0)
    {\iflangja{演算律速}{compute-bound}};
  % axes
  \draw[->] (-3,-3) -- (7.3,-3);
  \draw[->] (-3,-3) -- (-3,5.6);
  \foreach \e/\l in {-3/{1/8},-2/{1/4},-1/{1/2},0/1,1/2,2/4,3/8,4/16,5/32,6/64,7/128} {
    \draw (\e,-3) -- ++(0,-0.15) node[below, inner sep=1pt] {\l};
  }
  \foreach \e/\l in {-3/{1/8},-1/{1/2},1/2,3/8,4/16,5/32} {
    \draw (-3,\e) -- ++(-0.12,0) node[left, inner sep=1pt] {\l};
  }
  \node[below, inner sep=1pt] at (2,-4.1)
    {\iflangja{演算強度（FLOP/B，対数目盛）}{arithmetic intensity (FLOP/B, log scale)}};
  \node[rotate=90, above, inner sep=1pt] at (-4.15,1)
    {\iflangja{性能（TFLOP/s，対数目盛）}{performance (TFLOP/s, log scale)}};
  % roofs
  \draw[very thick, ln-accent] (-3,-3) -- (4,4);
  \draw[very thick, ln-tip] (4,4) -- (7,4);
  \node[above, text=ln-tip] at (5.5,4)
    {\iflangja{ピーク 16\,TFLOP/s}{peak 16\,TFLOP/s}};
  \node[rotate=35.5, below, text=ln-accent] at (2.7,2.7)
    {\iflangja{帯域 1\,TB/s $\times$ 演算強度}{1\,TB/s $\times$ intensity}};
  % ridge point
  \draw[dashed, ln-muted] (4,-3) -- (4,4);
  \fill[ln-caution] (4,4) circle (0.09);
  \node[right, text=ln-caution, align=left] at (4.1,3.25)
    {\iflangja{リッジポイント\\16\,FLOP/B}{ridge point\\16\,FLOP/B}};
  % kernels
  \fill (-1,-1) circle (0.1) node[below right, inner sep=1.5pt] {A: 0.5};
  \fill (1,1) circle (0.1) node[below right, inner sep=1.5pt] {B: 2};
  \draw[->, thick, ln-caution] (-0.85,-0.75) to[bend left=20] (0.8,0.85);
\end{tikzpicture}

% Width: body (120 mm); the drawing is about 110 mm wide.
\begin{tikzpicture}[x=8.5mm, y=6mm, font=\footnotesize\sffamily,
    >={Stealth[length=1.5mm]}]
  % regions (x = log2 FLOP/B from -3 to 7, y = log2 TFLOP/s from -3 to 5)
  \fill[ln-accent!8] (-3,-3) rectangle (4,5.4);
  \fill[ln-tip!10] (4,-3) rectangle (7,5.4);
  \node[text=ln-accent, align=center] at (-0.2,4.6)
    {\iflangja{メモリ律速}{memory-bound}};
  \node[text=ln-tip, align=center] at (5.5,0)
    {\iflangja{演算律速}{compute-bound}};
  % axes
  \draw[->] (-3,-3) -- (7.3,-3);
  \draw[->] (-3,-3) -- (-3,5.6);
  \foreach \e/\l in {-3/{1/8},-2/{1/4},-1/{1/2},0/1,1/2,2/4,3/8,4/16,5/32,6/64,7/128} {
    \draw (\e,-3) -- ++(0,-0.15) node[below, inner sep=1pt] {\l};
  }
  \foreach \e/\l in {-3/{1/8},-1/{1/2},1/2,3/8,4/16,5/32} {
    \draw (-3,\e) -- ++(-0.12,0) node[left, inner sep=1pt] {\l};
  }
  \node[below, inner sep=1pt] at (2,-4.1)
    {\iflangja{演算強度（FLOP/B，対数目盛）}{arithmetic intensity (FLOP/B, log scale)}};
  \node[rotate=90, above, inner sep=1pt] at (-4.15,1)
    {\iflangja{性能（TFLOP/s，対数目盛）}{performance (TFLOP/s, log scale)}};
  % roofs
  \draw[very thick, ln-accent] (-3,-3) -- (4,4);
  \draw[very thick, ln-tip] (4,4) -- (7,4);
  \node[above, text=ln-tip] at (5.5,4)
    {\iflangja{ピーク 16\,TFLOP/s}{peak 16\,TFLOP/s}};
  \node[rotate=35.5, below, text=ln-accent] at (2.7,2.7)
    {\iflangja{帯域 1\,TB/s $\times$ 演算強度}{1\,TB/s $\times$ intensity}};
  % ridge point
  \draw[dashed, ln-muted] (4,-3) -- (4,4);
  \fill[ln-caution] (4,4) circle (0.09);
  \node[right, text=ln-caution, align=left] at (4.1,3.25)
    {\iflangja{リッジポイント\\16\,FLOP/B}{ridge point\\16\,FLOP/B}};
  % kernels
  \fill (-1,-1) circle (0.1) node[below right, inner sep=1.5pt] {A: 0.5};
  \fill (1,1) circle (0.1) node[below right, inner sep=1.5pt] {B: 2};
  \draw[->, thick, ln-caution] (-0.85,-0.75) to[bend left=20] (0.8,0.85);
\end{tikzpicture}

% Width: body (120 mm); the drawing is about 110 mm wide.
\begin{tikzpicture}[x=8.5mm, y=6mm, font=\footnotesize\sffamily,
    >={Stealth[length=1.5mm]}]
  % regions (x = log2 FLOP/B from -3 to 7, y = log2 TFLOP/s from -3 to 5)
  \fill[ln-accent!8] (-3,-3) rectangle (4,5.4);
  \fill[ln-tip!10] (4,-3) rectangle (7,5.4);
  \node[text=ln-accent, align=center] at (-0.2,4.6)
    {\iflangja{メモリ律速}{memory-bound}};
  \node[text=ln-tip, align=center] at (5.5,0)
    {\iflangja{演算律速}{compute-bound}};
  % axes
  \draw[->] (-3,-3) -- (7.3,-3);
  \draw[->] (-3,-3) -- (-3,5.6);
  \foreach \e/\l in {-3/{1/8},-2/{1/4},-1/{1/2},0/1,1/2,2/4,3/8,4/16,5/32,6/64,7/128} {
    \draw (\e,-3) -- ++(0,-0.15) node[below, inner sep=1pt] {\l};
  }
  \foreach \e/\l in {-3/{1/8},-1/{1/2},1/2,3/8,4/16,5/32} {
    \draw (-3,\e) -- ++(-0.12,0) node[left, inner sep=1pt] {\l};
  }
  \node[below, inner sep=1pt] at (2,-4.1)
    {\iflangja{演算強度（FLOP/B，対数目盛）}{arithmetic intensity (FLOP/B, log scale)}};
  \node[rotate=90, above, inner sep=1pt] at (-4.15,1)
    {\iflangja{性能（TFLOP/s，対数目盛）}{performance (TFLOP/s, log scale)}};
  % roofs
  \draw[very thick, ln-accent] (-3,-3) -- (4,4);
  \draw[very thick, ln-tip] (4,4) -- (7,4);
  \node[above, text=ln-tip] at (5.5,4)
    {\iflangja{ピーク 16\,TFLOP/s}{peak 16\,TFLOP/s}};
  \node[rotate=35.5, below, text=ln-accent] at (2.7,2.7)
    {\iflangja{帯域 1\,TB/s $\times$ 演算強度}{1\,TB/s $\times$ intensity}};
  % ridge point
  \draw[dashed, ln-muted] (4,-3) -- (4,4);
  \fill[ln-caution] (4,4) circle (0.09);
  \node[right, text=ln-caution, align=left] at (4.1,3.25)
    {\iflangja{リッジポイント\\16\,FLOP/B}{ridge point\\16\,FLOP/B}};
  % kernels
  \fill (-1,-1) circle (0.1) node[below right, inner sep=1.5pt] {A: 0.5};
  \fill (1,1) circle (0.1) node[below right, inner sep=1.5pt] {B: 2};
  \draw[->, thick, ln-caution] (-0.85,-0.75) to[bend left=20] (0.8,0.85);
\end{tikzpicture}

  \caption{\SI{16}{\tera\flop\per\second} と
    \SI{1}{\tera\byte\per\second} のマシンのルーフラインモデル
    （リッジポイントは \SI{16}{\flop\per\byte}）．カーネル A
    （\SI{0.5}{\flop\per\byte}）は最大でも
    \SI{0.5}{\tera\flop\per\second} しか達成できない．演算強度を
    \SI{2}{\flop\per\byte} に上げると（カーネル B），上界は4倍になる．}
  \label{fig:l08-roofline}
\end{figure}

このモデルを導入した Williams らは，演算強度をキャッシュの後で測る．
彼らの\term[えんざんきょうどどらむばいとあたり]{演算強度（DRAMバイト当たり）}[operational intensity]
は DRAM トラフィック1バイトあたりの演算数を数えるので，データが
キャッシュにとどまるカーネルでは，operational intensity が
arithmetic intensity より高くなる
\cite{williams2009roofline}．\caution{講義 vs 論文：本章は講義と同様に，
プログラムが要求するバイト数を数える．論文は\enquote{the term
operational intensity to mean operations per byte of DRAM traffic}を
用い，上界を\enquote{Attainable GFlops/sec = Min(Peak Floating Point
Performance, Peak Memory Bandwidth x Operational Intensity)}と述べて
いる \cite{williams2009roofline}．}

一次近似設計の単位，すなわち要素あたり1バイトでの2つの例が，カーネルが
チャート上をどう動くかを示す．ベクトル和
\mbox{\lstinline[language=C++]!z[i] = x[i] + y[i]!} は，2つのロード
（それぞれ1バイト）だけを数えると2~バイトあたり1~FLOP を行う．その
演算強度は \SI{0.5}{\flop\per\byte} で，カーネルはメモリ律速である．
内積の変形
\mbox{\lstinline[language=C++]!sum += x[i]*x[i]*y[i]*y[i]!} は同じ
2~バイトを読むが，4~FLOP（3回の乗算と1回の加算）を行い，演算強度は
\SI{2}{\flop\per\byte} になる．チャート上で右へ，演算律速の領域に向かって
動く．このモデルの実用上の用途は，カーネルに改善の余地があるかを判断
することである．屋根のはるか下にあるカーネルは別の何か（レイテンシ，
ダイバージェンス，バンク競合）に制限されており，屋根の上にあるカーネルは
演算強度を上げるか，別のマシンに移ることでしか改善できない．屋根は
マシンごとに1回描かれ，\enquote{not once per kernel}である
\cite{williams2009roofline}．各カーネルは演算強度軸上の1点に
なる．\margin{実際の GPU のリッジポイント（FP16 テンソルのピークを HBM
帯域で割ったもの）：V100 \SI{138.9}{\flop\per\byte}
\cite{nvidia2026gemmguide}；A100 $312/1.555 \approx 201$；H100 SXM5
$1000/3 \approx 333$（暫定値 \cite{nvidia2020ampere,nvidia2022hopper}）．}

\begin{example}
  \label{ex:l08-roofline}
  \cref{fig:l08-roofline}の我々のマシンは
  $P_{\text{peak}} = \SI{16}{\tera\flop\per\second}$，
  $B_{\text{peak}} = \SI{1}{\tera\byte\per\second}$ なので，
  $I^{*} = 16/1 = \SI{16}{\flop\per\byte}$ である．FP32 の内積は FMA
  （2~FLOP）あたり \SI{8}{\byte} を読むので
  $I = 2/8 = \SI{0.25}{\flop\per\byte}$ であり，達成できるのは最大でも
  $1 \times 0.25 = \SI{0.25}{\tera\flop\per\second}$，ピークの約
  1.6\,\% である．上の4~FLOP の変形は FP32 で
  \SI{0.5}{\flop\per\byte} に達する．大きなタイル化行列積のように
  $I = \SI{32}{\flop\per\byte}$ のカーネルは
  $\min(16, 32) = \SI{16}{\tera\flop\per\second}$，すなわちピーク全体を
  達成する．
\end{example}

\subsection{CPI モデル}

解析モデルの第2の系統は，命令あたりの時間から出発する．

\begin{definition}[CPI]
  \label{def:l08-cpi}
  プロセッサ上のプログラムの
  \term[めいれいあたりさいくるすう]{命令当たりサイクル数（CPI）}[cycles per instruction (CPI)]
  は，それにかかるサイクル数を実行する命令数で割ったものであり，
  第1章の IPC の逆数である．
\end{definition}

最も単純な CPI モデルは，理想的なパイプラインの CPI に，あらゆる種類の
乱れの事象のコストを，その発生頻度で重み付けして加える．
\begin{equation}
  \mathrm{CPI} = \mathrm{CPI}_{\text{steady}} +
    \sum_{i} f_i \cdot p_i ,
  \label{eq:l08-cpi}
\end{equation}
ここで $f_i$ は命令あたりの種類 $i$ の事象の数，$p_i$ はその事象1回の
ペナルティ（サイクル）である．このモデルはペナルティが重ならない，
すなわち直列化されると仮定する．\caution{直列化された和は，ペナルティを
互いに，また有用な仕事と重ね合わせるアウトオブオーダープロセッサや
マルチスレッドプロセッサに対しては悲観的である．}

\begin{example}
  \label{ex:l08-cpi}
  5段のインオーダーパイプラインでは $\mathrm{CPI}_{\text{steady}} = 1$
  である．分岐予測ミスのコストは\num{3}~サイクルで命令の2\,\%で発生し，
  キャッシュミスのコストは\num{5}~サイクルで5\,\%で発生するとする．
  \[
    \mathrm{CPI} = 1 + 0.02 \cdot 3 + 0.05 \cdot 5 = 1 + 0.06 + 0.25
    = 1.31 .
  \]
\end{example}

\begin{marginfigure}
  \resizebox{\linewidth}{!}{% A CPI stack (Lesson 8): the CPI of 1.31 of the five-stage example
% split into its base (steady-state) component of 1.00 and the
% contributions of branch mispredictions (0.06) and cache misses (0.25).
% Original.
% Usage: \resizebox{\linewidth}{!}{% A CPI stack (Lesson 8): the CPI of 1.31 of the five-stage example
% split into its base (steady-state) component of 1.00 and the
% contributions of branch mispredictions (0.06) and cache misses (0.25).
% Original.
% Usage: \resizebox{\linewidth}{!}{% A CPI stack (Lesson 8): the CPI of 1.31 of the five-stage example
% split into its base (steady-state) component of 1.00 and the
% contributions of branch mispredictions (0.06) and cache misses (0.25).
% Original.
% Usage: \resizebox{\linewidth}{!}{\input{figures/l08-cpi-stack}}
% Width: margin (50 mm); the drawing is about 45 mm wide.
\begin{tikzpicture}[x=1mm, y=24mm, font=\scriptsize\sffamily]
  \draw[fill=ln-accent!20] (0,0) rectangle (12,1.00)
    node[midway, align=center] {\iflangja{基本}{base}\\1.00};
  \draw[fill=ln-tip!30] (0,1.00) rectangle (12,1.06);
  \draw[fill=ln-caution!25] (0,1.06) rectangle (12,1.31);
  \draw[ln-muted] (12,1.03) -- (15,1.03)
    node[right, text=black] {\iflangja{分岐予測ミス 0.06}{branch mispred.\ 0.06}};
  \draw[ln-muted] (12,1.185) -- (15,1.185)
    node[right, text=black] {\iflangja{キャッシュミス 0.25}{cache misses 0.25}};
  \draw[->] (-2,0) -- (-2,1.45) node[above, inner sep=1pt] {CPI};
  \foreach \v in {0,0.5,1} {
    \draw (-2,\v) -- (-2.8,\v) node[left, inner sep=1pt] {\v};
  }
  \draw[dashed, ln-muted] (-2,1.31) -- (0,1.31);
  \node[left, inner sep=1pt] at (-2.8,1.31) {1.31};
\end{tikzpicture}
}
% Width: margin (50 mm); the drawing is about 45 mm wide.
\begin{tikzpicture}[x=1mm, y=24mm, font=\scriptsize\sffamily]
  \draw[fill=ln-accent!20] (0,0) rectangle (12,1.00)
    node[midway, align=center] {\iflangja{基本}{base}\\1.00};
  \draw[fill=ln-tip!30] (0,1.00) rectangle (12,1.06);
  \draw[fill=ln-caution!25] (0,1.06) rectangle (12,1.31);
  \draw[ln-muted] (12,1.03) -- (15,1.03)
    node[right, text=black] {\iflangja{分岐予測ミス 0.06}{branch mispred.\ 0.06}};
  \draw[ln-muted] (12,1.185) -- (15,1.185)
    node[right, text=black] {\iflangja{キャッシュミス 0.25}{cache misses 0.25}};
  \draw[->] (-2,0) -- (-2,1.45) node[above, inner sep=1pt] {CPI};
  \foreach \v in {0,0.5,1} {
    \draw (-2,\v) -- (-2.8,\v) node[left, inner sep=1pt] {\v};
  }
  \draw[dashed, ln-muted] (-2,1.31) -- (0,1.31);
  \node[left, inner sep=1pt] at (-2.8,1.31) {1.31};
\end{tikzpicture}
}
% Width: margin (50 mm); the drawing is about 45 mm wide.
\begin{tikzpicture}[x=1mm, y=24mm, font=\scriptsize\sffamily]
  \draw[fill=ln-accent!20] (0,0) rectangle (12,1.00)
    node[midway, align=center] {\iflangja{基本}{base}\\1.00};
  \draw[fill=ln-tip!30] (0,1.00) rectangle (12,1.06);
  \draw[fill=ln-caution!25] (0,1.06) rectangle (12,1.31);
  \draw[ln-muted] (12,1.03) -- (15,1.03)
    node[right, text=black] {\iflangja{分岐予測ミス 0.06}{branch mispred.\ 0.06}};
  \draw[ln-muted] (12,1.185) -- (15,1.185)
    node[right, text=black] {\iflangja{キャッシュミス 0.25}{cache misses 0.25}};
  \draw[->] (-2,0) -- (-2,1.45) node[above, inner sep=1pt] {CPI};
  \foreach \v in {0,0.5,1} {
    \draw (-2,\v) -- (-2.8,\v) node[left, inner sep=1pt] {\v};
  }
  \draw[dashed, ln-muted] (-2,1.31) -- (0,1.31);
  \node[left, inner sep=1pt] at (-2.8,1.31) {1.31};
\end{tikzpicture}
}
  \caption{\cref{ex:l08-cpi}の CPI スタック．}
  \label{fig:l08-cpi-stack}
\end{marginfigure}

和の各項は別々に示す価値がある．時間がどこに使われているかを
示すからである．

\begin{definition}[CPI スタック]
  \label{def:l08-cpi-stack}
  \term[しーぴーあいすたっく]{CPIスタック}[CPI stack]は，CPI の成分を
  積み上げた棒グラフである．基本（定常状態）の成分と，事象の種類
  ごとに1つの成分からなる（\cref{fig:l08-cpi-stack}）．
\end{definition}

この例では，キャッシュミスが時間の $0.25/1.31 \approx 19\,\%$ を，
予測ミスが5\,\% 未満を占める．よりよい予測器よりよいキャッシュの方が
効果が大きい．\tip{IPC は CPI の逆数である．CPI 1.31 は IPC
$1/1.31 \approx 0.76$ である．}

マルチスレッドプロセッサは，あるスレッドのストールを他のスレッドの
仕事と重ね合わせる．理想的には，$W_{\text{depth}}$ 個のスレッドで
ストールはそれらの間で分割される．
\begin{equation}
  \mathrm{CPI}_{\text{MT}} \approx
    \frac{\mathrm{CPI}_{\text{single}}}{W_{\text{depth}}}
    + \text{contention terms}.
  \label{eq:l08-cpi-mt}
\end{equation}
競合項（contention terms）は，多くのスレッドが動いているときに
足りなくなる共有資源を表す．ビジーな実行ユニット，DRAM 帯域，そして
未処理のミスを追跡するレジスタである．

\begin{definition}[MSHR]
  \label{def:l08-mshr}
  \term{MSHR}[MSHR (miss status holding register)]（miss status
  holding register）は，未処理のキャッシュミス1つ（ミスしたブロックと，
  それを待つ要求）を記録するエントリであり，これによってキャッシュは
  ミスの処理中にも他のアクセスに応じ続けられる．MSHR の数が，同時に
  処理中のミスの数を制限する．
\end{definition}
\margin{したがって MSHR は第4章のメモリレベル並列性に上限を課す．
たとえば32個の MSHR（説明のための数値）では，準備のできたワープが
いくつあっても，33個目の未処理ミスは1つが戻るまでストールする．}

\begin{note}
  \label{note:l08-cpi-mt}
  \eqref{eq:l08-cpi-mt}は，すべてのスレッドの命令を数えた，プロセッサ
  全体の命令あたりサイクル数として読む．理想項は発行の限界（1サイクルに
  1命令を発行するパイプラインでは1）を下回れない．\cref{ex:l08-cpi}の
  CPI 1.31 では，2スレッドですでにその限界に達する．GPU は多数の
  ワープを持つマルチスレッドプロセッサなので CPI モデルが当てはまるが，
  競合項，特に MSHR と DRAM 帯域が支配的になる．
\end{note}

\subsection{インターバル解析}

CPI モデルは，各ペナルティを，それ以外は完璧なパイプラインに加わる固定
サイクル数として扱う．インターバル解析は，実行の時間的な形に目を向ける．

\begin{definition}[インターバル解析]
  \label{def:l08-interval}
  \term[いんたーばるかいせき]{インターバル解析}[interval analysis]の
  手法は，実行をインターバルに分割する．各インターバルは，命令が
  \emph{実効ディスパッチ率}でディスパッチされる定常状態と，それに
  続く1つの\term[みすいべんと]{ミスイベント}[miss event]からなる．
  ミスイベントは長いレイテンシの事象（分岐予測ミス，キャッシュミスや
  TLB ミス）であり，そのペナルティの間は命令がディスパッチされない
  （\cref{fig:l08-interval}）．実行時間は定常状態の時間とペナルティの
  和である．
\end{definition}

\begin{figure}[tb]
  \centering
  % Interval analysis (Lesson 8): dispatch rate over time.  Each interval
% is a steady state at the effective dispatch rate followed by one miss
% event.  Dashed: the simple model, in which dispatch stops abruptly for
% the penalty.  Solid: the first-order model, with the pipeline drain
% before and the refill after the penalty.
% Redraws the outline figure m8l3_interval_based (both panels merged).
% Usage: % Interval analysis (Lesson 8): dispatch rate over time.  Each interval
% is a steady state at the effective dispatch rate followed by one miss
% event.  Dashed: the simple model, in which dispatch stops abruptly for
% the penalty.  Solid: the first-order model, with the pipeline drain
% before and the refill after the penalty.
% Redraws the outline figure m8l3_interval_based (both panels merged).
% Usage: % Interval analysis (Lesson 8): dispatch rate over time.  Each interval
% is a steady state at the effective dispatch rate followed by one miss
% event.  Dashed: the simple model, in which dispatch stops abruptly for
% the penalty.  Solid: the first-order model, with the pipeline drain
% before and the refill after the penalty.
% Redraws the outline figure m8l3_interval_based (both panels merged).
% Usage: \input{figures/l08-interval}
% Width: body (120 mm); the drawing is about 112 mm wide.
\begin{tikzpicture}[x=1mm, y=1mm, font=\scriptsize\sffamily,
    >={Stealth[length=1.5mm]}]
  \draw[->] (0,0) -- (110,0) node[below left, inner sep=1pt]
    {\iflangja{時間}{time}};
  \draw[->] (0,0) -- (0,24);
  \node[rotate=90, align=center, inner sep=1pt] at (-7,8)
    {\iflangja{ディスパッチ率}{dispatch rate}};
  \node[left, inner sep=1pt] at (0,14) {$D$};
  % simple model (dashed)
  \draw[dashed, ln-muted] (2,14) -- (36,14) -- (36,0.3) -- (56,0.3)
    -- (56,14) -- (90,14) -- (90,0.3) -- (104,0.3);
  % first-order model (solid)
  \draw[very thick, ln-accent] (2,14) -- (32,14) -- (38,0.6) -- (50,0.6)
    -- (58,14) -- (86,14) -- (92,0.6) -- (104,0.6);
  % labels
  \draw[->, ln-caution] (32,21) node[above, text=ln-caution, inner sep=1pt]
    {\iflangja{ミスイベント}{miss event}} -- (32,15);
  \draw[->, ln-caution] (86,21) node[above, text=ln-caution, inner sep=1pt]
    {\iflangja{ミスイベント}{miss event}} -- (86,15);
  \node[above, text=ln-accent] at (17,14)
    {\iflangja{定常状態}{steady state}};
  \node[text=ln-muted] at (35,6) [left] {\iflangja{ドレイン}{drain}};
  \node[text=ln-muted] at (55,6) [right] {\iflangja{リフィル}{refill}};
  \node[above, inner sep=1pt] at (45,1) {\iflangja{ペナルティ}{penalty}};
  % intervals
  \draw[|-|] (2,-4) -- (56,-4) node[midway, below, inner sep=1.5pt]
    {\iflangja{インターバル 1}{interval 1}};
  \draw[|-|] (56,-4) -- (104,-4) node[midway, below, inner sep=1.5pt]
    {\iflangja{インターバル 2}{interval 2}};
\end{tikzpicture}

% Width: body (120 mm); the drawing is about 112 mm wide.
\begin{tikzpicture}[x=1mm, y=1mm, font=\scriptsize\sffamily,
    >={Stealth[length=1.5mm]}]
  \draw[->] (0,0) -- (110,0) node[below left, inner sep=1pt]
    {\iflangja{時間}{time}};
  \draw[->] (0,0) -- (0,24);
  \node[rotate=90, align=center, inner sep=1pt] at (-7,8)
    {\iflangja{ディスパッチ率}{dispatch rate}};
  \node[left, inner sep=1pt] at (0,14) {$D$};
  % simple model (dashed)
  \draw[dashed, ln-muted] (2,14) -- (36,14) -- (36,0.3) -- (56,0.3)
    -- (56,14) -- (90,14) -- (90,0.3) -- (104,0.3);
  % first-order model (solid)
  \draw[very thick, ln-accent] (2,14) -- (32,14) -- (38,0.6) -- (50,0.6)
    -- (58,14) -- (86,14) -- (92,0.6) -- (104,0.6);
  % labels
  \draw[->, ln-caution] (32,21) node[above, text=ln-caution, inner sep=1pt]
    {\iflangja{ミスイベント}{miss event}} -- (32,15);
  \draw[->, ln-caution] (86,21) node[above, text=ln-caution, inner sep=1pt]
    {\iflangja{ミスイベント}{miss event}} -- (86,15);
  \node[above, text=ln-accent] at (17,14)
    {\iflangja{定常状態}{steady state}};
  \node[text=ln-muted] at (35,6) [left] {\iflangja{ドレイン}{drain}};
  \node[text=ln-muted] at (55,6) [right] {\iflangja{リフィル}{refill}};
  \node[above, inner sep=1pt] at (45,1) {\iflangja{ペナルティ}{penalty}};
  % intervals
  \draw[|-|] (2,-4) -- (56,-4) node[midway, below, inner sep=1.5pt]
    {\iflangja{インターバル 1}{interval 1}};
  \draw[|-|] (56,-4) -- (104,-4) node[midway, below, inner sep=1.5pt]
    {\iflangja{インターバル 2}{interval 2}};
\end{tikzpicture}

% Width: body (120 mm); the drawing is about 112 mm wide.
\begin{tikzpicture}[x=1mm, y=1mm, font=\scriptsize\sffamily,
    >={Stealth[length=1.5mm]}]
  \draw[->] (0,0) -- (110,0) node[below left, inner sep=1pt]
    {\iflangja{時間}{time}};
  \draw[->] (0,0) -- (0,24);
  \node[rotate=90, align=center, inner sep=1pt] at (-7,8)
    {\iflangja{ディスパッチ率}{dispatch rate}};
  \node[left, inner sep=1pt] at (0,14) {$D$};
  % simple model (dashed)
  \draw[dashed, ln-muted] (2,14) -- (36,14) -- (36,0.3) -- (56,0.3)
    -- (56,14) -- (90,14) -- (90,0.3) -- (104,0.3);
  % first-order model (solid)
  \draw[very thick, ln-accent] (2,14) -- (32,14) -- (38,0.6) -- (50,0.6)
    -- (58,14) -- (86,14) -- (92,0.6) -- (104,0.6);
  % labels
  \draw[->, ln-caution] (32,21) node[above, text=ln-caution, inner sep=1pt]
    {\iflangja{ミスイベント}{miss event}} -- (32,15);
  \draw[->, ln-caution] (86,21) node[above, text=ln-caution, inner sep=1pt]
    {\iflangja{ミスイベント}{miss event}} -- (86,15);
  \node[above, text=ln-accent] at (17,14)
    {\iflangja{定常状態}{steady state}};
  \node[text=ln-muted] at (35,6) [left] {\iflangja{ドレイン}{drain}};
  \node[text=ln-muted] at (55,6) [right] {\iflangja{リフィル}{refill}};
  \node[above, inner sep=1pt] at (45,1) {\iflangja{ペナルティ}{penalty}};
  % intervals
  \draw[|-|] (2,-4) -- (56,-4) node[midway, below, inner sep=1.5pt]
    {\iflangja{インターバル 1}{interval 1}};
  \draw[|-|] (56,-4) -- (104,-4) node[midway, below, inner sep=1.5pt]
    {\iflangja{インターバル 2}{interval 2}};
\end{tikzpicture}

  \caption{インターバル解析．各インターバルは実効ディスパッチ率 $D$ の
    定常状態と，それに続く1つのミスイベントからなる．破線：単純な
    モデルの急停止．実線：ペナルティの前のパイプラインのドレインと
    後のリフィルを含む一次近似モデル．}
  \label{fig:l08-interval}
\end{figure}

パイプラインは急に止まって急に再開するわけではない．ミスイベントが
起こっても，すでに処理中の命令はパイプラインから流れ出ていく．
それが解決した後も，ディスパッチが定常の速さに達するまでには
パイプラインが再び満たされる時間が必要である．Karkhanis と Smith の
一次近似スーパースカラモデルはこの\emph{ドレイン}と\emph{リフィル}の
段階を加える．その形は事象の種類による．予測ミスは分岐が実行されて
初めて検出され，長いキャッシュミスはリオーダバッファの先頭を塞ぎ，
バッファが満杯になるとディスパッチを止める
\cite{karkhanis2004firstorder}．\margin{Eeckhout
\cite{eeckhout2010perfeval} はインターバル解析と一次近似モデルを，
ミスイベントの種類ごとのペナルティも含めて詳しく述べている．}

GPU に適用するとき，素朴な方法は GPU をもう1つのマルチスレッド
プロセッサとして扱う．それでは GPU を最もよく特徴づける2つの効果を
取り逃がす．
\begin{description}
  \item[分岐ダイバージェンス．]  ワープのすべてのレーンがアクティブとは
    限らず，ダイバージェント分岐の経路は直列化される（第6章）ので，
    ワープ命令は一定量の仕事を表さない．
  \item[メモリダイバージェンス．]  ワープのメモリ命令のレイテンシは，
    そのアクセスがコアレッシングされるかどうかに依存する．
    コアレッシングされないアクセスは多数のトランザクションを生成し，
    その最後のものを待つ（第4章）．
\end{description}
\margin{GPUMech はインターバル解析を GPU に適用し，第7章のラウンド
ロビンと greedy-then-oldest のスケジューラを，シミュレーションの97倍の
速さで13.2\,\% と14.0\,\% の誤差でモデル化し，CPI スタックを生成する
\cite{huang2014gpumech}．}%
より新しい GPU のインターバルモデルは，サブコア，セクタキャッシュ，
その他の共有資源の競合を加えている．

\begin{note}[メモリと計算の並列性による GPU モデル]
  \label{note:l08-mwp-cwp}
  本ユニットの任意の読み物の1つである Hong と Kim の GPU の解析モデルは，
  2つの量に基づく \cite{hong2009analytical}．
  \term[めもりわーぷへいれつど]{メモリワープ並列度}[memory warp parallelism (MWP)]
  （MWP）は，1つのワープのメモリ待ち期間中に SM あたり同時にメモリに
  アクセスできるワープの最大数である．それはレイテンシ，連続する要求の
  間の発行間隔（departure delay），帯域，アクティブなワープ数 $N$ で
  制限される．
  $\mathrm{MWP} = \min(\mathrm{MWP}_{\text{without BW}},
  \mathrm{MWP}_{\text{peak BW}}, N)$ であり，ここで
  $\mathrm{MWP}_{\text{without BW}} =
  \textit{Mem\_L}/\textit{Departure\_delay}$，
  $\mathrm{MWP}_{\text{peak BW}}$ はピーク帯域を，1ワープが必要とする
  帯域とアクティブな SM の数の積で割ったものである．
  \term[けいさんわーぷへいれつど]{計算ワープ並列度}[computation warp parallelism (CWP)]
  （CWP）は，1回のメモリ待ち期間中に SM が実行できるワープの数に1を
  加えたものである．
  $\mathrm{CWP} = \min\bigl((\textit{Mem\_cycles} +
  \textit{Comp\_cycles})/\textit{Comp\_cycles},\; N\bigr)$．
  $\mathrm{CWP} \geq \mathrm{MWP}$ ならば，他のワープの計算は
  メモリシステムが重ね合わせられる以上を覆う．カーネルはメモリ律速で
  あり，その時間は $N/\mathrm{MWP}$ 回のメモリ期間に比例して増える．
  $\mathrm{MWP} > \mathrm{CWP}$ ならメモリレイテンシは隠され，時間は
  およそメモリレイテンシ1回分と $N \times \textit{Comp\_cycles}$ の
  和である．
\end{note}
\margin{Hong と Kim は，マイクロベンチマークで5.4\,\%，アプリケーションで
13.3\,\% の幾何平均誤差を報告している \cite{hong2009analytical}．}

\begin{keypoint}
  解析モデルは精度と引き換えに洞察を得る．ルーフラインはカーネルの
  演算強度をマシンのリッジポイントと比べ，演算律速かメモリ律速かを
  判定する．CPI モデルは定常状態の CPI に頻度で重み付けしたペナルティを
  加え，CPI スタックは時間がどこに使われるかを示す．インターバル解析は
  これをドレインとリフィルで精密にし，GPU 版は分岐ダイバージェンスと
  メモリダイバージェンスを加えなければならない．
\end{keypoint}

\section{GPU シミュレーションの高速化}
\label{sec:l08-accelerating}

詳細なシミュレーションは遅い．\source{M8，シミュレーションの高速化；
Eeckhout \cite{eeckhout2010perfeval}（サンプリング）；\cite{liu2023photon,sherwood2002simpoint}．}
各サイクルに各ユニットをモデル化するサイクルレベルシミュレータは，
1秒あたりシミュレーション上の命令千個程度（\num{1} KIPS）の速さで動く．
その速さでは，プロセッサが約1秒で実行する\num{e9}個の命令の
シミュレーションに約28~時間かかり，実ハードウェアで数時間動く機械学習の
ワークロードには1世紀程度かかる．\caution{講義 vs 計算：
$\num{e9}/\num{e3} = \num{e6}$~s は約11.6~日であり，28~時間ではない．
28~時間は10~KIPS（$\num{e5}$~s）に相当する．1~KIPS では，ネイティブ
実行1時間分に約 $3600 \times 11.6~\text{日} \approx 114$~年かかり，
これが1世紀である．}出口は2つある．シミュレータを速くするか，
シミュレーションする量を減らすかである．

\subsection{より速いシミュレータ}

\paragraph{並列シミュレーション．}  シミュレータは複数のホストスレッド，
たとえばシミュレーションするコアごとに1スレッドで動かせる．\margin{CPU
シミュレータの速さ：gem5，Flexus，MARSS は詳細なモデルで約200~KIPS に
達する．zsim は詳細なアウトオブオーダーコアを20~MIPS で，1024コアの
チップを16コアのホスト上で最大300~MIPS でシミュレーションする
\cite{sanchez2013zsim}．}コアはほぼ独立しているが，メモリシステムと
ネットワークオンチップは共有されている．スレッドは相互作用のたびに
通信するか，共有モデルへのアクセスを直列化しなければならない．
ほとんどのワークロードはメモリ律速なので，こうした相互作用は頻繁で，
メモリモデルが並列シミュレータのボトルネックになる．

\paragraph{イベント駆動シミュレーション．}

\begin{definition}[イベント駆動シミュレーション]
  \label{def:l08-event-driven}
  \term[いべんとくどうしみゅれーしょん]{イベント駆動シミュレーション}[event-driven simulation]
  では，シミュレータはタイムスタンプ付きの将来のイベントのキューを
  時刻順に保持する．最も早いイベントを取り出し，シミュレーション時刻を
  そのタイムスタンプまで直接進めて処理することを繰り返す．処理は新しい
  イベントを挿入しうる（\cref{fig:l08-event-queue}）．
\end{definition}

\begin{figure}[tb]
  \centering
  % Event-driven simulation (Lesson 8).  Top: an event queue ordered by
% time stamp; processing the event at cycle 100 inserts a new event at
% cycle 250 in time order.  Bottom: simulated time jumps from one event
% to the next, skipping the cycles in between.
% Redraws the outline figure m8l4_esim with our own events and times.
% Usage: % Event-driven simulation (Lesson 8).  Top: an event queue ordered by
% time stamp; processing the event at cycle 100 inserts a new event at
% cycle 250 in time order.  Bottom: simulated time jumps from one event
% to the next, skipping the cycles in between.
% Redraws the outline figure m8l4_esim with our own events and times.
% Usage: % Event-driven simulation (Lesson 8).  Top: an event queue ordered by
% time stamp; processing the event at cycle 100 inserts a new event at
% cycle 250 in time order.  Bottom: simulated time jumps from one event
% to the next, skipping the cycles in between.
% Redraws the outline figure m8l4_esim with our own events and times.
% Usage: \input{figures/l08-event-queue}
% Width: body (120 mm); the drawing is about 112 mm wide.
\begin{tikzpicture}[x=1mm, y=1mm, font=\scriptsize\sffamily,
    >={Stealth[length=1.5mm]},
    ev/.style={draw, fill=ln-box, minimum width=24mm, minimum height=8mm,
               align=center, inner sep=1pt}]
  \node[anchor=west, font=\scriptsize\sffamily\bfseries] at (0,40)
    {\iflangja{イベントキュー（時刻順）}{event queue (ordered by time)}};
  \node[ev, fill=ln-accent!15] (e1) at (14,31)
    {$t=100$\\\iflangja{DRAM 応答}{DRAM reply}};
  \node[ev] (e2) at (42,31)
    {$t=104$\\\iflangja{ワープ3 起床}{warp 3 wakes}};
  \node[ev] (e3) at (98,31)
    {$t=400$\\\iflangja{I/O 完了}{I/O done}};
  \node[ev, fill=ln-tip!15] (e4) at (70,31)
    {$t=250$\\\iflangja{L2 書き戻し}{L2 write-back}};
  \draw[->] (e1) -- (e2);
  \draw[->] (e2) -- (e4);
  \draw[->] (e4) -- (e3);
  \draw[->, ln-tip] (e1.south) to[bend right=18]
    node[pos=0.55, below, text=ln-tip, inner sep=1.5pt]
    {\iflangja{処理中に新しいイベントを挿入}{processing it inserts a new event}}
    (e4.south);
  % timeline
  \draw[->] (0,6) -- (110,6);
  \node[anchor=north east, inner sep=1pt] at (110,0.5)
    {\iflangja{シミュレーション時刻（サイクル）}{simulated time (cycles)}};
  \foreach \t/\x in {100/14,104/22,250/62,400/98} {
    \fill (\x,6) circle (0.8);
    \node[below, inner sep=2pt] at (\x,5) {\t};
  }
  \draw[->, ln-caution] (14,7.5) to[bend left=50] (22,7.5);
  \draw[->, ln-caution] (22,7.5) to[bend left=30]
    node[midway, above, text=ln-caution, inner sep=1pt]
    {\iflangja{145サイクルをスキップ}{skip 145 cycles}} (62,7.5);
  \draw[->, ln-caution] (62,7.5) to[bend left=30]
    node[midway, above, text=ln-caution, inner sep=1pt]
    {\iflangja{149サイクルをスキップ}{skip 149 cycles}} (98,7.5);
\end{tikzpicture}

% Width: body (120 mm); the drawing is about 112 mm wide.
\begin{tikzpicture}[x=1mm, y=1mm, font=\scriptsize\sffamily,
    >={Stealth[length=1.5mm]},
    ev/.style={draw, fill=ln-box, minimum width=24mm, minimum height=8mm,
               align=center, inner sep=1pt}]
  \node[anchor=west, font=\scriptsize\sffamily\bfseries] at (0,40)
    {\iflangja{イベントキュー（時刻順）}{event queue (ordered by time)}};
  \node[ev, fill=ln-accent!15] (e1) at (14,31)
    {$t=100$\\\iflangja{DRAM 応答}{DRAM reply}};
  \node[ev] (e2) at (42,31)
    {$t=104$\\\iflangja{ワープ3 起床}{warp 3 wakes}};
  \node[ev] (e3) at (98,31)
    {$t=400$\\\iflangja{I/O 完了}{I/O done}};
  \node[ev, fill=ln-tip!15] (e4) at (70,31)
    {$t=250$\\\iflangja{L2 書き戻し}{L2 write-back}};
  \draw[->] (e1) -- (e2);
  \draw[->] (e2) -- (e4);
  \draw[->] (e4) -- (e3);
  \draw[->, ln-tip] (e1.south) to[bend right=18]
    node[pos=0.55, below, text=ln-tip, inner sep=1.5pt]
    {\iflangja{処理中に新しいイベントを挿入}{processing it inserts a new event}}
    (e4.south);
  % timeline
  \draw[->] (0,6) -- (110,6);
  \node[anchor=north east, inner sep=1pt] at (110,0.5)
    {\iflangja{シミュレーション時刻（サイクル）}{simulated time (cycles)}};
  \foreach \t/\x in {100/14,104/22,250/62,400/98} {
    \fill (\x,6) circle (0.8);
    \node[below, inner sep=2pt] at (\x,5) {\t};
  }
  \draw[->, ln-caution] (14,7.5) to[bend left=50] (22,7.5);
  \draw[->, ln-caution] (22,7.5) to[bend left=30]
    node[midway, above, text=ln-caution, inner sep=1pt]
    {\iflangja{145サイクルをスキップ}{skip 145 cycles}} (62,7.5);
  \draw[->, ln-caution] (62,7.5) to[bend left=30]
    node[midway, above, text=ln-caution, inner sep=1pt]
    {\iflangja{149サイクルをスキップ}{skip 149 cycles}} (98,7.5);
\end{tikzpicture}

% Width: body (120 mm); the drawing is about 112 mm wide.
\begin{tikzpicture}[x=1mm, y=1mm, font=\scriptsize\sffamily,
    >={Stealth[length=1.5mm]},
    ev/.style={draw, fill=ln-box, minimum width=24mm, minimum height=8mm,
               align=center, inner sep=1pt}]
  \node[anchor=west, font=\scriptsize\sffamily\bfseries] at (0,40)
    {\iflangja{イベントキュー（時刻順）}{event queue (ordered by time)}};
  \node[ev, fill=ln-accent!15] (e1) at (14,31)
    {$t=100$\\\iflangja{DRAM 応答}{DRAM reply}};
  \node[ev] (e2) at (42,31)
    {$t=104$\\\iflangja{ワープ3 起床}{warp 3 wakes}};
  \node[ev] (e3) at (98,31)
    {$t=400$\\\iflangja{I/O 完了}{I/O done}};
  \node[ev, fill=ln-tip!15] (e4) at (70,31)
    {$t=250$\\\iflangja{L2 書き戻し}{L2 write-back}};
  \draw[->] (e1) -- (e2);
  \draw[->] (e2) -- (e4);
  \draw[->] (e4) -- (e3);
  \draw[->, ln-tip] (e1.south) to[bend right=18]
    node[pos=0.55, below, text=ln-tip, inner sep=1.5pt]
    {\iflangja{処理中に新しいイベントを挿入}{processing it inserts a new event}}
    (e4.south);
  % timeline
  \draw[->] (0,6) -- (110,6);
  \node[anchor=north east, inner sep=1pt] at (110,0.5)
    {\iflangja{シミュレーション時刻（サイクル）}{simulated time (cycles)}};
  \foreach \t/\x in {100/14,104/22,250/62,400/98} {
    \fill (\x,6) circle (0.8);
    \node[below, inner sep=2pt] at (\x,5) {\t};
  }
  \draw[->, ln-caution] (14,7.5) to[bend left=50] (22,7.5);
  \draw[->, ln-caution] (22,7.5) to[bend left=30]
    node[midway, above, text=ln-caution, inner sep=1pt]
    {\iflangja{145サイクルをスキップ}{skip 145 cycles}} (62,7.5);
  \draw[->, ln-caution] (62,7.5) to[bend left=30]
    node[midway, above, text=ln-caution, inner sep=1pt]
    {\iflangja{149サイクルをスキップ}{skip 149 cycles}} (98,7.5);
\end{tikzpicture}

  \caption{イベント駆動シミュレーション．イベントキューはタイムスタンプ順に
    並び，イベントの処理は新しいイベントを挿入しうる．シミュレーション
    時刻はイベントからイベントへ跳び，その間のサイクルを飛ばす．}
  \label{fig:l08-event-queue}
\end{figure}

イベント間のサイクルにはコストがかからないので，イベント駆動
シミュレーションはイベントがまばらなとき，たとえばレイテンシが長いときに
効果を発揮する．数千サイクルかかる I/O 操作や不揮発性メモリへの書き込みは，
数千回のティックではなく1つのイベントになる．\tip{サイクルレベル
シミュレーションのコストはサイクルごとに，イベント駆動シミュレーションの
コストはイベントごとにかかる．忙しいパイプラインはほぼ毎サイクルに
イベントがあるのでコアの得るものは小さく，レイテンシの長いデバイスが
最も得をする．}

\paragraph{より単純なモデル．}  すべてのイベントが同じように重要なわけでは
ない．シミュレータは重要でないイベントには平均のレイテンシや
スループットを使い，問いが必要とするところだけ詳細を保つことができる．
メモリ階層を調べるなら，コアを発行幅に等しい IPC で命令を発行する
だけのものに縮め，キャッシュとメモリだけをモデル化すればよい．
多コア CPU の高速シミュレータである zsim は，詳細なアウトオブオーダー
モデルと並んで，そのような単純なコア（1サイクルに1命令）を提供する
\cite{sanchez2013zsim}．コアを調べるなら，メモリに固定のレイテンシを
与えればよい．\caution{どの単純化も相互作用を1つ取り除く．たとえば固定の
メモリレイテンシは，帯域を大量に消費する設計が生むキューイング遅延，
まさに\eqref{eq:l08-cpi-mt}の競合項を隠してしまう．}

残る2つの高速化，サンプリングと統計モデルが，本節の残りの主題である．

\subsection{サンプリング}

\begin{definition}[サンプリングシミュレーション]
  \label{def:l08-sampled}
  \term[さんぷりんぐしみゅれーしょん]{サンプリングシミュレーション}[sampled simulation]
  では，プログラムの実行のうち選んだ部分（サンプル）だけを詳細に
  シミュレーションし，それらから全体の性能を推定する．サンプルは
  ランダムに選んでも，プログラムを解析して選んでもよい．
\end{definition}

サンプルの間の部分も何らかの方法で飛ばさなければならない．実行駆動
シミュレータは，タイミングなしで機能的に実行する（はるかに速い）ことで
そこを早送りするか，各サンプルを
\term[ちぇっくぽいんと]{チェックポイント}[checkpoint]から開始する．
チェックポイントは保存されたアーキテクチャ状態（レジスタとメモリ）で
あり，そこからそれ以前の部分を実行せずにシミュレーションを再開できる．
トレース駆動シミュレータは，単にサンプルのトレースだけを収集する．

サンプルは現実的なマイクロアーキテクチャ状態から始まるわけではない．
そのキャッシュと分岐予測器は空であるか，前のサンプルが残した内容を
保持している．サンプルの\term[うぉーむあっぷ]{ウォームアップ}[warm-up]は，
その前にシミュレーションされ，タイミングは捨てられる実行区間であり，
その間にこれらの構造が満たされる．あるいは，コールドスタートが無視
できるほどサンプルを長くする．

ランダムなサンプルで正確な推定を得るには多くの抽出が必要である．しかし
プログラムはランダムではなく，同じ振る舞いを繰り返す．

\begin{definition}[プログラムフェーズと SimPoint]
  \label{def:l08-phase}
  \term[ぷろぐらむふぇーず]{プログラムフェーズ}[program phase]は，
  振る舞い（IPC，ミス率）が均質な実行区間であり，何度も繰り返されうる．
  \term{SimPoint}[SimPoint]は，プログラムの実行のインターバルを
  それが実行するコードによってクラスタリングすることでプログラムの
  フェーズを識別し，フェーズごとに1つの代表インターバルを，そのフェーズが
  実行に占める割合で重み付けして選ぶ（\cref{fig:l08-phases}）
  \cite{sherwood2002simpoint}．各インターバルは\emph{基本ブロック
  ベクトル}，すなわちその中で各基本ブロックが何回実行されたかで記述され，
  ベクトルは $k$-means クラスタリングでまとめられる．
\end{definition}

\begin{marginfigure}
  \resizebox{\linewidth}{!}{% Program phases and phase-based sampling (Lesson 8): IPC over time for
% a program that alternates between a high-IPC phase A and a low-IPC
% phase B.  Phase-based sampling (SimPoint) simulates one representative
% interval per phase (outlined) and weights it by the phase's share.
% Redraws the outline figure m8l4_sample with our own phase lengths.
% Usage: \resizebox{\linewidth}{!}{% Program phases and phase-based sampling (Lesson 8): IPC over time for
% a program that alternates between a high-IPC phase A and a low-IPC
% phase B.  Phase-based sampling (SimPoint) simulates one representative
% interval per phase (outlined) and weights it by the phase's share.
% Redraws the outline figure m8l4_sample with our own phase lengths.
% Usage: \resizebox{\linewidth}{!}{% Program phases and phase-based sampling (Lesson 8): IPC over time for
% a program that alternates between a high-IPC phase A and a low-IPC
% phase B.  Phase-based sampling (SimPoint) simulates one representative
% interval per phase (outlined) and weights it by the phase's share.
% Redraws the outline figure m8l4_sample with our own phase lengths.
% Usage: \resizebox{\linewidth}{!}{\input{figures/l08-phases}}
% Width: margin (50 mm); the drawing is about 50 mm wide.
\begin{tikzpicture}[x=1.15mm, y=1mm, font=\scriptsize\sffamily]
  \foreach \a/\b/\p/\h/\c in {0/8/A/16/ln-tip!30,8/14/B/6/ln-source!25,
      14/22/A/16/ln-tip!30,22/28/B/6/ln-source!25,28/36/A/16/ln-tip!30,
      36/40/B/6/ln-source!25} {
    \fill[\c] (\a,0) rectangle (\b,\h);
    \node at ({(\a+\b)/2},-2.2) {\p};
  }
  \draw[very thick, ln-caution] (2,0) rectangle (5,16);
  \draw[very thick, ln-caution] (9.5,0) rectangle (12.5,6);
  \draw[->] (0,0) -- (41.5,0);
  \draw[->] (0,0) -- (0,19) node[above right, inner sep=1pt] {IPC};
  \node[anchor=north east] at (41.5,-3.8) {\iflangja{時間}{time}};
\end{tikzpicture}
}
% Width: margin (50 mm); the drawing is about 50 mm wide.
\begin{tikzpicture}[x=1.15mm, y=1mm, font=\scriptsize\sffamily]
  \foreach \a/\b/\p/\h/\c in {0/8/A/16/ln-tip!30,8/14/B/6/ln-source!25,
      14/22/A/16/ln-tip!30,22/28/B/6/ln-source!25,28/36/A/16/ln-tip!30,
      36/40/B/6/ln-source!25} {
    \fill[\c] (\a,0) rectangle (\b,\h);
    \node at ({(\a+\b)/2},-2.2) {\p};
  }
  \draw[very thick, ln-caution] (2,0) rectangle (5,16);
  \draw[very thick, ln-caution] (9.5,0) rectangle (12.5,6);
  \draw[->] (0,0) -- (41.5,0);
  \draw[->] (0,0) -- (0,19) node[above right, inner sep=1pt] {IPC};
  \node[anchor=north east] at (41.5,-3.8) {\iflangja{時間}{time}};
\end{tikzpicture}
}
% Width: margin (50 mm); the drawing is about 50 mm wide.
\begin{tikzpicture}[x=1.15mm, y=1mm, font=\scriptsize\sffamily]
  \foreach \a/\b/\p/\h/\c in {0/8/A/16/ln-tip!30,8/14/B/6/ln-source!25,
      14/22/A/16/ln-tip!30,22/28/B/6/ln-source!25,28/36/A/16/ln-tip!30,
      36/40/B/6/ln-source!25} {
    \fill[\c] (\a,0) rectangle (\b,\h);
    \node at ({(\a+\b)/2},-2.2) {\p};
  }
  \draw[very thick, ln-caution] (2,0) rectangle (5,16);
  \draw[very thick, ln-caution] (9.5,0) rectangle (12.5,6);
  \draw[->] (0,0) -- (41.5,0);
  \draw[->] (0,0) -- (0,19) node[above right, inner sep=1pt] {IPC};
  \node[anchor=north east] at (41.5,-3.8) {\iflangja{時間}{time}};
\end{tikzpicture}
}
  \caption{IPC の高いフェーズ A と低いフェーズ B を交互に繰り返す
    プログラムの IPC の時間変化．フェーズに基づくサンプリングは，
    それぞれの代表1つ（枠で囲んだもの）をシミュレーションする．}
  \label{fig:l08-phases}
\end{marginfigure}

少数のフェーズが数百回繰り返されるなら，各フェーズの1インターバルを
シミュレーションして重み付けすることで，完全なシミュレーションの精度の
大部分を，そのコストのわずかな割合で得られる．

\subsection{GPU でのサンプリング}

GPU の1スレッドは短すぎてその内部でサンプリングすることはできないが，
GPU のプログラムは同じコードを実行する非常に多くのスレッドからなり
（SPMD，第1章），別の粒度が使える．
\begin{description}
  \item[ブロック単位．]  グリッド全体の数千のスレッドブロックの代わりに，
    100個程度のスレッドブロックをシミュレーションする．
  \item[カーネル単位．]  数十のカーネルの代わりに1つか2つのカーネルを
    シミュレーションする．
  \item[ワープ単位．]  カーネルのより少ないワープを詳細にシミュレーション
    する．
\end{description}

\begin{definition}[カーネル単位サンプリング]
  \label{def:l08-kernel-sampling}
  \term[かーねるたんいさんぷりんぐ]{カーネル単位サンプリング}[kernel-level sampling]
  は，異なるカーネルごとに，あるいは類似した起動のグループごとに1回の
  起動を詳細にシミュレーションし，その時間を他の起動に用いる．
\end{definition}

カーネル単位サンプリングは，層と反復ごとに1回ずつ同じカーネルを数百回
起動する機械学習のワークロードで特に効果的である．本ユニットの指定の
読み物の1つである Photon は，カーネル，ワープ，基本ブロックの3つの
レベルのサンプリングを組み合わせ，事前にワークロードをプロファイルする
代わりに軽量なオンライン解析によって反復的な振る舞いを見つける
\cite{liu2023photon}．\margin{Photon は ResNet-152 の1回の推論
（バッチサイズ~1）のシミュレーションを7.05~日から1.7~時間に短縮し，
サンプリング誤差は10.7\,\% である．MGPUSim の上に構築されている
\cite{liu2023photon}．}

\subsection{ワークロードの縮小}

もう1つの出口は，シミュレーションする仕事を減らすことである．
\begin{itemize}
  \item \emph{反復を減らす}：機械学習の訓練の反復は似た振る舞いをする
    ので，1000回ではなく1回の訓練反復をシミュレーションする．
  \item \emph{入力を小さくする}：\num{e9}ノードではなく\num{e6}ノードの
    グラフ．
  \item \emph{ホットなコードだけ}：数百の関数のうち1つか2つが時間の
    90\,\% を占めるなら，それらだけをシミュレーションする．
\end{itemize}
\caution{小さな入力は振る舞いを短くするだけでなく変えることがある．
縮小したサイズでは L2 キャッシュに収まるワーキングセットが完全な
サイズでは収まらないことがあり，縮小した実行は実際の実行の
メモリ律速な振る舞いを隠しうる．}

\begin{definition}[マイクロベンチマークと代表ワークロード]
  \label{def:l08-microbenchmark}
  \term[まいくろべんちまーく]{マイクロベンチマーク}[microbenchmark]は，
  マシンの1つの機能（あるキャッシュレベル，あるメモリアクセスパターン，
  ある命令のレイテンシ）を単独で動かすために書かれた小さなプログラムで
  ある．\term[だいひょうわーくろーど]{代表ワークロード}[representative workload]
  は，調べる機能について振る舞いが完全なワークロードと一致するように
  作った縮小版のプログラムまたは入力である．
\end{definition}

\subsection{データ駆動モデル}

最後の方法はマシンをまったくモデル化しない．

\begin{definition}[データ駆動モデル]
  \label{def:l08-data-driven}
  \term[でーたくどうもでる]{データ駆動モデル}[data-driven model]は，
  ハードウェアの動作を組み込むことなく，データに当てはめた統計モデルや
  機械学習モデルによって，プログラムの特徴量（命令数，ロード数，
  ダイバージェント分岐の数，\dots）から性能を予測する．
\end{definition}

\Needspace{20\baselineskip}
手順は3段階である．一連のプログラムをハードウェア（または詳細な
シミュレータ）で実行してその性能，たとえば IPC を記録する．ハードウェア
性能カウンタまたはシミュレータから特徴量を抽出する．そしてモデルを
当てはめる．\margin{ハードウェアカウンタは Nsight Compute（第5章）が
読むものである．種類ごとの実行命令数，メモリトランザクション，分岐
ダイバージェンスなど．}そのようなモデルの単純な形は
\begin{equation}
  \textit{cycles} = C_1 N_{\text{inst}}^{e_1}
    + C_2 N_{\text{FP}}^{e_2} + C_3 N_{\text{mem}}^{e_3}
  \label{eq:l08-data-driven}
\end{equation}
である．ここで $N_{\text{inst}}$，$N_{\text{FP}}$，$N_{\text{mem}}$ は
全命令，浮動小数点命令，メモリ命令の数であり，係数 $C_k$ と指数
$e_k$ はデータに当てはめて決める．\caution{解析モデルはマシンの理解から
作られ，データ駆動モデルはデータだけから作られる．それは当てはめに
使ったものと似たプログラムに対しては正確だが，その外では何の保証も
与えず，洞察もほとんど与えない．}

\begin{keypoint}[章のまとめ]
  サイクルレベルシミュレータはクロックを進めて各サイクルに各ユニットを
  モデル化し，実行またはトレースによって駆動される．パイプラインは深さと
  幅を持つ FIFO キューの集合とスケジューラであり，GPU シミュレータは
  ワープ，そのマスクまたは SIMT スタック，コアレッシング，セクタ
  キャッシュをモデル化する．解析モデルは一次近似の答えを与える．設計の
  FMA ユニット，帯域，$W_{\text{depth}} \approx L/n$，リッジポイントを
  持つルーフラインの上界 $\min(P_{\text{peak}}, B_{\text{peak}} \cdot I)$，
  CPI スタック，分岐ダイバージェンスとメモリダイバージェンスを含む
  インターバル解析である．約1~KIPS のシミュレーションは実際の
  ワークロードには遅すぎる．並列シミュレーションとイベント駆動
  シミュレーション，単純化したモデル，サンプリング（フェーズ，SimPoint，
  ブロック単位，カーネル単位，ワープ単位），縮小したワークロード，
  データ駆動モデルによって，それを扱える規模にする．
\end{keypoint}

\end{document}
